\section{未锚定账本事件的叙事补充：three-kingdoms}

\label{sec:anchor-batch-two-three-kingdoms}

以下事件均为此前正文尚未设置导航目标的账本记录。每一锚点置于来源、制度与地方条件的叙事中，避免把年表、政权分类或战报修辞当作社会全貌。

\subsection{制度与地方资源}

\eventanchor{TK-0253-ZHUGE-KE-HEFEI}{建兴二年（253年）·诸葛恪围攻合肥新城久攻不下，吴军因疫病、补给与魏守备而撤退}账本记录孙吴与曹魏在建兴二年（253年）发生“诸葛恪围攻合肥新城久攻不下，吴军因疫病、补给与魏守备而撤退”。条目把背景概括为东兴胜利提高诸葛恪威望，他试图趁势夺取合肥以改变淮南战略格局，却低估了攻城补给与夏季疾病成本，并指出远征损害吴军和诸葛恪的政治信誉，为其被孙峻等人推翻创造条件；疾病规模不能由史书修辞量化。这个节点进入区域社会时，要经由文书、道路、军队、市场、寺院和家庭等不同中介；因此，政治行动的宣布、地方接收和日常生活的改变不能被视作同一时刻完成。

本条依据《三国志·诸葛恪传》；《三国志·孙亮传》；《三国志·毌丘俭传》；《资治通鉴》魏纪，证据提示为“围城、撤军和军中疫病可靠；兵数、围城日期和败因权重有争议”。这些材料可支持年序、官方决定、战事或特定地点的线索，却难以直接统计所有人的财富、迁徙、信仰和动机。更稳妥的读法是区分诏令与实施、军报与人口、行政称谓与自我认同，并让地方材料的缺环限制解释的范围。

由此可见，该事件不是孤立的年表点，而是资源、信息与权力关系重新配置的一环。不同地区的官员、社群和家庭可能以不同方式承担征发、保护、迁移或宗教活动；现存来源只能照亮其中一部分，不能替沉默者制造统一的经验。

\eventanchor{TK-0253-SUN-JUN-COUP}{建兴二年（253年）·孙峻在建业杀诸葛恪并接掌禁军与朝政}账本记录孙吴在建兴二年（253年）发生“孙峻在建业杀诸葛恪并接掌禁军与朝政”。条目把背景概括为诸葛恪的合肥失败激化朝臣与宗室对其专权的恐惧，孙峻控制宫廷武力并利用这些不满，并指出孙吴辅政由外戚重臣转到宗室将领孙峻手中，孙亮的皇权更加依赖禁军统帅而非公开朝议。这个节点进入区域社会时，要经由文书、道路、军队、市场、寺院和家庭等不同中介；因此，政治行动的宣布、地方接收和日常生活的改变不能被视作同一时刻完成。

本条依据《三国志·诸葛恪传》；《三国志·孙峻传》；《三国志·孙亮传》；《资治通鉴》魏纪，证据提示为“政变与诸葛恪死亡可靠；孙亮知情程度、宫廷程序和孙峻动机有争议”。这些材料可支持年序、官方决定、战事或特定地点的线索，却难以直接统计所有人的财富、迁徙、信仰和动机。更稳妥的读法是区分诏令与实施、军报与人口、行政称谓与自我认同，并让地方材料的缺环限制解释的范围。

由此可见，该事件不是孤立的年表点，而是资源、信息与权力关系重新配置的一环。不同地区的官员、社群和家庭可能以不同方式承担征发、保护、迁移或宗教活动；现存来源只能照亮其中一部分，不能替沉默者制造统一的经验。

\eventanchor{TK-0254-CAO-FANG-DEPOSED}{正元元年（254年）·司马师废曹芳，改立高贵乡公曹髦为帝}账本记录曹魏在正元元年（254年）发生“司马师废曹芳，改立高贵乡公曹髦为帝”。条目把背景概括为司马师与曹芳身边近臣关系恶化，司马氏已掌握禁军、太后名义和高层官僚的关键节点，并指出曹魏皇帝可由权臣更换成为公开先例，曹髦虽具文辞和政治意志，却继承了极为有限的实权。这个节点进入区域社会时，要经由文书、道路、军队、市场、寺院和家庭等不同中介；因此，政治行动的宣布、地方接收和日常生活的改变不能被视作同一时刻完成。

本条依据《三国志·齐王芳纪》；《三国志·高贵乡公纪》；《三国志·司马师传》；《资治通鉴》魏纪，证据提示为“废立与继位可靠；曹芳被废罪名、太后诏令的形成和曹髦被选中的过程有争议”。这些材料可支持年序、官方决定、战事或特定地点的线索，却难以直接统计所有人的财富、迁徙、信仰和动机。更稳妥的读法是区分诏令与实施、军报与人口、行政称谓与自我认同，并让地方材料的缺环限制解释的范围。

由此可见，该事件不是孤立的年表点，而是资源、信息与权力关系重新配置的一环。不同地区的官员、社群和家庭可能以不同方式承担征发、保护、迁移或宗教活动；现存来源只能照亮其中一部分，不能替沉默者制造统一的经验。

\eventanchor{TK-0255-GUANQIU-WEN-REBELLION}{正元二年（255年）·毌丘俭、文钦在淮南起兵反司马师，失败后文钦等投向孙吴}账本记录曹魏、淮南叛军与孙吴在正元二年（255年）发生“毌丘俭、文钦在淮南起兵反司马师，失败后文钦等投向孙吴”。条目把背景概括为曹芳被废使魏室支持者担忧司马氏将取代曹氏，淮南将领拥有接近吴境的兵员与退路，并指出司马氏平定淮南而孙吴获得流亡将领，魏吴边界成为反司马氏政治与跨境军事互动的通道。这个节点进入区域社会时，要经由文书、道路、军队、市场、寺院和家庭等不同中介；因此，政治行动的宣布、地方接收和日常生活的改变不能被视作同一时刻完成。

本条依据《三国志·毌丘俭传》；《三国志·文钦传》；《三国志·司马师传》；《资治通鉴》魏纪，证据提示为“起兵、失败和文钦南奔可靠；叛军规模、孙吴事前支持和淮南民众反应有争议”。这些材料可支持年序、官方决定、战事或特定地点的线索，却难以直接统计所有人的财富、迁徙、信仰和动机。更稳妥的读法是区分诏令与实施、军报与人口、行政称谓与自我认同，并让地方材料的缺环限制解释的范围。

由此可见，该事件不是孤立的年表点，而是资源、信息与权力关系重新配置的一环。不同地区的官员、社群和家庭可能以不同方式承担征发、保护、迁移或宗教活动；现存来源只能照亮其中一部分，不能替沉默者制造统一的经验。

\eventanchor{TK-0255-SIMA-SHI-DEATH}{正元二年闰正月（255年）·司马师在镇压淮南叛乱后病逝，司马昭接掌辅政与中军权力}账本记录曹魏在正元二年闰正月（255年）发生“司马师在镇压淮南叛乱后病逝，司马昭接掌辅政与中军权力”。条目把背景概括为司马师虽平定叛乱，却因病无法长期留在前线，司马氏必须在军事危机中维持兄弟间的权力连续，并指出司马昭成为实际执政者，继续面对曹髦、淮南和吴蜀的压力；权力交接并非意味着魏的反对力量已消失。这个节点进入区域社会时，要经由文书、道路、军队、市场、寺院和家庭等不同中介；因此，政治行动的宣布、地方接收和日常生活的改变不能被视作同一时刻完成。

本条依据《三国志·司马师传》；《三国志·司马昭传》；《三国志·高贵乡公纪》，证据提示为“卒年与司马昭继位可靠；病因、交接时的朝臣选择和司马昭权力基础有争议”。这些材料可支持年序、官方决定、战事或特定地点的线索，却难以直接统计所有人的财富、迁徙、信仰和动机。更稳妥的读法是区分诏令与实施、军报与人口、行政称谓与自我认同，并让地方材料的缺环限制解释的范围。

由此可见，该事件不是孤立的年表点，而是资源、信息与权力关系重新配置的一环。不同地区的官员、社群和家庭可能以不同方式承担征发、保护、迁移或宗教活动；现存来源只能照亮其中一部分，不能替沉默者制造统一的经验。

\eventanchor{TK-0256-SUN-JUN-DEATH}{太平元年（256年）·孙峻去世，孙綝继掌禁军和朝政，孙亮试图摆脱权臣的空间更小}账本记录孙吴在太平元年（256年）发生“孙峻去世，孙綝继掌禁军和朝政，孙亮试图摆脱权臣的空间更小”。条目把背景概括为孙峻以宗室身份控制幼主政权后未建立可受朝廷监督的辅政机制，其死亡使军权向同族孙綝转移，并指出孙綝继续支配诏令和宫禁，孙吴的继承问题与权臣控制更深地交织；宗室称谓不应掩盖具体军政网络。这个节点进入区域社会时，要经由文书、道路、军队、市场、寺院和家庭等不同中介；因此，政治行动的宣布、地方接收和日常生活的改变不能被视作同一时刻完成。

本条依据《三国志·孙峻传》；《三国志·孙綝传》；《三国志·孙亮传》，证据提示为“孙峻死亡和孙綝继权可靠；孙峻死因、孙綝接班程序和孙亮参与程度有争议”。这些材料可支持年序、官方决定、战事或特定地点的线索，却难以直接统计所有人的财富、迁徙、信仰和动机。更稳妥的读法是区分诏令与实施、军报与人口、行政称谓与自我认同，并让地方材料的缺环限制解释的范围。

由此可见，该事件不是孤立的年表点，而是资源、信息与权力关系重新配置的一环。不同地区的官员、社群和家庭可能以不同方式承担征发、保护、迁移或宗教活动；现存来源只能照亮其中一部分，不能替沉默者制造统一的经验。

\eventanchor{TK-0257-ZHUGE-DAN-REBELLION}{甘露二年（257年）·诸葛诞据寿春反司马昭，并向孙吴求援，淮南爆发第三次大规模反司马氏叛乱}账本记录曹魏、淮南叛军与孙吴在甘露二年（257年）发生“诸葛诞据寿春反司马昭，并向孙吴求援，淮南爆发第三次大规模反司马氏叛乱”。条目把背景概括为司马氏先后废帝、平定淮南，使持有重兵的诸葛诞预期自身将被清除；寿春又可依淮河接受吴援，并指出司马昭被迫集中魏军围寿春，孙吴试图借叛乱扭转淮南均势；叛乱的长期供给成为决定性问题。这个节点进入区域社会时，要经由文书、道路、军队、市场、寺院和家庭等不同中介；因此，政治行动的宣布、地方接收和日常生活的改变不能被视作同一时刻完成。

本条依据《三国志·诸葛诞传》；《三国志·司马昭传》；《三国志·孙亮传》；《资治通鉴》魏纪，证据提示为“叛乱、求吴援助和寿春对峙可靠；兵员、粮储、吴军投入和诸葛诞动机有争议”。这些材料可支持年序、官方决定、战事或特定地点的线索，却难以直接统计所有人的财富、迁徙、信仰和动机。更稳妥的读法是区分诏令与实施、军报与人口、行政称谓与自我认同，并让地方材料的缺环限制解释的范围。

由此可见，该事件不是孤立的年表点，而是资源、信息与权力关系重新配置的一环。不同地区的官员、社群和家庭可能以不同方式承担征发、保护、迁移或宗教活动；现存来源只能照亮其中一部分，不能替沉默者制造统一的经验。

\eventanchor{TK-0258-ZHUGE-DAN-DEFEAT}{甘露三年（258年）·寿春陷落，诸葛诞被杀，司马昭击败吴援军并结束淮南连续叛乱}账本记录曹魏、淮南叛军与孙吴在甘露三年（258年）发生“寿春陷落，诸葛诞被杀，司马昭击败吴援军并结束淮南连续叛乱”。条目把背景概括为魏军以长期围困切断吴援与城内补给，叛军内部在持久战和降附问题上分裂，并指出司马昭清除淮南最后的强大反对中心，魏的军事资源得以转向蜀汉；寿春战后人口与处置规模须谨慎处理。这个节点进入区域社会时，要经由文书、道路、军队、市场、寺院和家庭等不同中介；因此，政治行动的宣布、地方接收和日常生活的改变不能被视作同一时刻完成。

本条依据《三国志·诸葛诞传》；《三国志·司马昭传》；《三国志·孙綝传》；《资治通鉴》魏纪，证据提示为“城破、诸葛诞死亡和叛乱终结可靠；城内饥困、投降过程、屠杀与吴军损失有争议”。这些材料可支持年序、官方决定、战事或特定地点的线索，却难以直接统计所有人的财富、迁徙、信仰和动机。更稳妥的读法是区分诏令与实施、军报与人口、行政称谓与自我认同，并让地方材料的缺环限制解释的范围。

由此可见，该事件不是孤立的年表点，而是资源、信息与权力关系重新配置的一环。不同地区的官员、社群和家庭可能以不同方式承担征发、保护、迁移或宗教活动；现存来源只能照亮其中一部分，不能替沉默者制造统一的经验。

\eventanchor{TK-0258-SUN-LIANG-DEPOSED}{太平三年（258年）·孙綝废孙亮，迎立琅邪王孙休，孙吴第三位皇帝即位}账本记录孙吴在太平三年（258年）发生“孙綝废孙亮，迎立琅邪王孙休，孙吴第三位皇帝即位”。条目把背景概括为孙亮试图联合近侍和宗室制衡孙綝但失败，孙綝以禁军优势操控宫廷并寻找更易接受的成年宗室，并指出孙吴皇帝再次由权臣更换，孙休随后必须设法处置孙綝；频繁废立削弱了吴对外动员的可信度。这个节点进入区域社会时，要经由文书、道路、军队、市场、寺院和家庭等不同中介；因此，政治行动的宣布、地方接收和日常生活的改变不能被视作同一时刻完成。

本条依据《三国志·孙亮传》；《三国志·孙休传》；《三国志·孙綝传》；《资治通鉴》魏纪，证据提示为“废立和孙休即位可靠；孙亮被废的指控、太后与宗室的作用及孙休受拥立的条件有争议”。这些材料可支持年序、官方决定、战事或特定地点的线索，却难以直接统计所有人的财富、迁徙、信仰和动机。更稳妥的读法是区分诏令与实施、军报与人口、行政称谓与自我认同，并让地方材料的缺环限制解释的范围。

由此可见，该事件不是孤立的年表点，而是资源、信息与权力关系重新配置的一环。不同地区的官员、社群和家庭可能以不同方式承担征发、保护、迁移或宗教活动；现存来源只能照亮其中一部分，不能替沉默者制造统一的经验。

\eventanchor{TK-0260-CAO-MAO-KILLED}{甘露五年（260年）·曹髦率近侍出宫试图讨伐司马昭，途中被成济杀害，司马昭另立曹奂}账本记录曹魏在甘露五年（260年）发生“曹髦率近侍出宫试图讨伐司马昭，途中被成济杀害，司马昭另立曹奂”。条目把背景概括为曹髦长期受司马昭监控而试图以皇帝名义动员近侍和官员，双方已无法在废立体制内共存，并指出皇帝在都城街道被杀使曹魏名义遭受不可逆损伤，曹奂即位后司马昭距离禅代仅剩礼仪和时机问题。这个节点进入区域社会时，要经由文书、道路、军队、市场、寺院和家庭等不同中介；因此，政治行动的宣布、地方接收和日常生活的改变不能被视作同一时刻完成。

本条依据《三国志·高贵乡公纪》；《三国志·司马昭传》；《三国志·王沈传》裴松之注；《资治通鉴》魏纪，证据提示为“曹髦出宫、被杀和曹奂继位可靠；曹髦的预谋、成济受命链条和司马昭即时责任有争议”。这些材料可支持年序、官方决定、战事或特定地点的线索，却难以直接统计所有人的财富、迁徙、信仰和动机。更稳妥的读法是区分诏令与实施、军报与人口、行政称谓与自我认同，并让地方材料的缺环限制解释的范围。

由此可见，该事件不是孤立的年表点，而是资源、信息与权力关系重新配置的一环。不同地区的官员、社群和家庭可能以不同方式承担征发、保护、迁移或宗教活动；现存来源只能照亮其中一部分，不能替沉默者制造统一的经验。

\eventanchor{TK-0262-JIANG-WEI-DIDAO}{景耀五年（262年）·姜维围攻狄道，邓艾据险救援，蜀军因粮运和魏军反击撤退}账本记录蜀汉与曹魏在景耀五年（262年）发生“姜维围攻狄道，邓艾据险救援，蜀军因粮运和魏军反击撤退”。条目把背景概括为姜维持续寻求在陇右打开缺口，蜀汉却需从汉中远运粮食，魏将领可依熟悉的山地道路增援，并指出蜀汉未能取得陇右据点并消耗有限军力，魏获得邓艾等熟悉蜀地边防的将领，为次年伐蜀积累情报与信心。这个节点进入区域社会时，要经由文书、道路、军队、市场、寺院和家庭等不同中介；因此，政治行动的宣布、地方接收和日常生活的改变不能被视作同一时刻完成。

本条依据《三国志·姜维传》；《三国志·邓艾传》；《三国志·后主传》；《资治通鉴》魏纪，证据提示为“围城、邓艾救援和撤军可靠；兵力、战果、郭淮角色与撤军损失有争议”。这些材料可支持年序、官方决定、战事或特定地点的线索，却难以直接统计所有人的财富、迁徙、信仰和动机。更稳妥的读法是区分诏令与实施、军报与人口、行政称谓与自我认同，并让地方材料的缺环限制解释的范围。

由此可见，该事件不是孤立的年表点，而是资源、信息与权力关系重新配置的一环。不同地区的官员、社群和家庭可能以不同方式承担征发、保护、迁移或宗教活动；现存来源只能照亮其中一部分，不能替沉默者制造统一的经验。

\eventanchor{TK-0263-WEI-INVADES-SHU}{景元四年（263年）·司马昭命钟会、邓艾、诸葛绪多路伐蜀，魏军由关中和陇右南下}账本记录曹魏与蜀汉在景元四年（263年）发生“司马昭命钟会、邓艾、诸葛绪多路伐蜀，魏军由关中和陇右南下”。条目把背景概括为淮南叛乱平定后司马昭可集中魏的关中军力，蜀汉的北伐消耗和朝廷内部矛盾又削弱了边防应变，并指出战争将决定三国均势是否还能维持；多路进军迫使姜维在剑阁、阴平和成都之间分散判断与兵力。这个节点进入区域社会时，要经由文书、道路、军队、市场、寺院和家庭等不同中介；因此，政治行动的宣布、地方接收和日常生活的改变不能被视作同一时刻完成。

本条依据《三国志·邓艾传》；《三国志·钟会传》；《三国志·姜维传》；《三国志·司马昭传》；《资治通鉴》魏纪，证据提示为“出师、主要统帅和多路部署可靠；兵数、诸葛绪失职、阴平路线筹划及各军进度有争议”。这些材料可支持年序、官方决定、战事或特定地点的线索，却难以直接统计所有人的财富、迁徙、信仰和动机。更稳妥的读法是区分诏令与实施、军报与人口、行政称谓与自我认同，并让地方材料的缺环限制解释的范围。

由此可见，该事件不是孤立的年表点，而是资源、信息与权力关系重新配置的一环。不同地区的官员、社群和家庭可能以不同方式承担征发、保护、迁移或宗教活动；现存来源只能照亮其中一部分，不能替沉默者制造统一的经验。

\eventanchor{TK-0263-CHENGDU-SURRENDER}{景耀六年（263年）·邓艾越阴平至成都平原，刘禅出降，蜀汉政权终结}账本记录蜀汉与曹魏在景耀六年（263年）发生“邓艾越阴平至成都平原，刘禅出降，蜀汉政权终结”。条目把背景概括为剑阁方向的姜维牵制钟会，邓艾的险路突进绕过主要防线，成都朝廷难以在短期内重建野战军，并指出三国之一被魏吞并，姜维等仍在剑阁而降诏已改变其行动空间；投降过程应避免用单一人物道德解释整个国家崩解。这个节点进入区域社会时，要经由文书、道路、军队、市场、寺院和家庭等不同中介；因此，政治行动的宣布、地方接收和日常生活的改变不能被视作同一时刻完成。

本条依据《三国志·后主传》；《三国志·邓艾传》；《三国志·谯周传》；《华阳国志·刘后主志》，证据提示为“邓艾入成都平原、刘禅降魏和政权终结可靠；阴平行军细节、谯周劝降作用、刘禅选择和成都守军状况有争议”。这些材料可支持年序、官方决定、战事或特定地点的线索，却难以直接统计所有人的财富、迁徙、信仰和动机。更稳妥的读法是区分诏令与实施、军报与人口、行政称谓与自我认同，并让地方材料的缺环限制解释的范围。

由此可见，该事件不是孤立的年表点，而是资源、信息与权力关系重新配置的一环。不同地区的官员、社群和家庭可能以不同方式承担征发、保护、迁移或宗教活动；现存来源只能照亮其中一部分，不能替沉默者制造统一的经验。

\eventanchor{TK-0263-SIMA-ZHAO-DUKE-JIN}{景元四年十二月（263年）·曹奂进封司马昭为晋公、加九锡，司马氏取得禅代前的封国框架}账本记录曹魏在景元四年十二月（263年）发生“曹奂进封司马昭为晋公、加九锡，司马氏取得禅代前的封国框架”。条目把背景概括为灭蜀大捷使司马昭的军功和朝廷控制达到高点，曹魏皇帝与官僚只能以传统封授语言确认既成权力，并指出晋公国为日后晋王和受禅建立官僚、礼仪与功臣安置框架；封爵文本是权力关系的表述，非平等协商的证据。这个节点进入区域社会时，要经由文书、道路、军队、市场、寺院和家庭等不同中介；因此，政治行动的宣布、地方接收和日常生活的改变不能被视作同一时刻完成。

本条依据《三国志·陈留王纪》；《三国志·司马昭传》；《资治通鉴》魏纪，证据提示为“进封和加九锡可靠；册命的自愿性、礼仪过程和晋公国实际机构范围有争议”。这些材料可支持年序、官方决定、战事或特定地点的线索，却难以直接统计所有人的财富、迁徙、信仰和动机。更稳妥的读法是区分诏令与实施、军报与人口、行政称谓与自我认同，并让地方材料的缺环限制解释的范围。

由此可见，该事件不是孤立的年表点，而是资源、信息与权力关系重新配置的一环。不同地区的官员、社群和家庭可能以不同方式承担征发、保护、迁移或宗教活动；现存来源只能照亮其中一部分，不能替沉默者制造统一的经验。

\eventanchor{TK-0264-ZHONG-HUI-REBELLION}{咸熙元年（264年）·钟会在成都图谋自立，拘杀邓艾后遭魏军与降众反杀，西南军政陷入短暂混乱}账本记录魏军、蜀地降众与钟会集团在咸熙元年（264年）发生“钟会在成都图谋自立，拘杀邓艾后遭魏军与降众反杀，西南军政陷入短暂混乱”。条目把背景概括为蜀汉骤亡后魏军统帅握有新占区域和大量降附军民，钟会又与邓艾围绕军功、任命和司马昭猜忌冲突，并指出司马昭虽保住西南，却失去两名重要将领并须重新安置蜀地军政；裴注异说须与陈寿正文分层使用。这个节点进入区域社会时，要经由文书、道路、军队、市场、寺院和家庭等不同中介；因此，政治行动的宣布、地方接收和日常生活的改变不能被视作同一时刻完成。

本条依据《三国志·钟会传》；《三国志·邓艾传》；《三国志·姜维传》；《三国志》裴松之注引《汉晋春秋》，证据提示为“钟会反叛、邓艾遇害和钟会死亡可靠；姜维参与、兵变细节、邓艾罪状及死者数量有争议”。这些材料可支持年序、官方决定、战事或特定地点的线索，却难以直接统计所有人的财富、迁徙、信仰和动机。更稳妥的读法是区分诏令与实施、军报与人口、行政称谓与自我认同，并让地方材料的缺环限制解释的范围。

由此可见，该事件不是孤立的年表点，而是资源、信息与权力关系重新配置的一环。不同地区的官员、社群和家庭可能以不同方式承担征发、保护、迁移或宗教活动；现存来源只能照亮其中一部分，不能替沉默者制造统一的经验。

\eventanchor{TK-0264-SIMA-ZHAO-KING-JIN}{咸熙元年三月（264年）·曹奂进封司马昭为晋王，封国等级进一步逼近受禅}账本记录曹魏在咸熙元年三月（264年）发生“曹奂进封司马昭为晋王，封国等级进一步逼近受禅”。条目把背景概括为司马昭灭蜀并平钟会后掌控魏的主要军队和新占土地，晋公名位已不足以承载其政治地位，并指出魏皇帝保留而司马氏王国扩张，禅代从可能性变为已制度化的预期；吴仍是阻止统一的唯一独立政权。这个节点进入区域社会时，要经由文书、道路、军队、市场、寺院和家庭等不同中介；因此，政治行动的宣布、地方接收和日常生活的改变不能被视作同一时刻完成。

本条依据《三国志·陈留王纪》；《三国志·司马昭传》；《晋书·文帝纪》，证据提示为“进封年月可靠；九锡、王国官署与朝臣劝进过程有争议”。这些材料可支持年序、官方决定、战事或特定地点的线索，却难以直接统计所有人的财富、迁徙、信仰和动机。更稳妥的读法是区分诏令与实施、军报与人口、行政称谓与自我认同，并让地方材料的缺环限制解释的范围。

由此可见，该事件不是孤立的年表点，而是资源、信息与权力关系重新配置的一环。不同地区的官员、社群和家庭可能以不同方式承担征发、保护、迁移或宗教活动；现存来源只能照亮其中一部分，不能替沉默者制造统一的经验。

\eventanchor{TK-0264-SUN-HAO-ACCESSION}{元兴元年（264年）·孙休去世后，孙皓被拥立为帝，孙吴进入末代皇帝统治}账本记录孙吴在元兴元年（264年）发生“孙休去世后，孙皓被拥立为帝，孙吴进入末代皇帝统治”。条目把背景概括为孙休未能留下稳固成年继承人，孙吴在孙綝被诛后仍需以宗室继承重建宫廷秩序，并指出孙皓继位时魏已灭蜀、晋王将受禅，孙吴面对更孤立的国际环境；后世暴君叙事须与具体政令和地方材料区分。这个节点进入区域社会时，要经由文书、道路、军队、市场、寺院和家庭等不同中介；因此，政治行动的宣布、地方接收和日常生活的改变不能被视作同一时刻完成。

本条依据《三国志·孙皓传》；《三国志·孙休传》；《建康实录》；《资治通鉴》魏纪，证据提示为“孙休卒与孙皓即位可靠；孙皓早期施政评价、拥立者动机和宗室选立过程有争议”。这些材料可支持年序、官方决定、战事或特定地点的线索，却难以直接统计所有人的财富、迁徙、信仰和动机。更稳妥的读法是区分诏令与实施、军报与人口、行政称谓与自我认同，并让地方材料的缺环限制解释的范围。

由此可见，该事件不是孤立的年表点，而是资源、信息与权力关系重新配置的一环。不同地区的官员、社群和家庭可能以不同方式承担征发、保护、迁移或宗教活动；现存来源只能照亮其中一部分，不能替沉默者制造统一的经验。

\eventanchor{TK-0265-SIMA-ZHAO-DEATH}{咸熙二年八月（265年）·司马昭去世，司马炎继承晋王位及魏廷实际军政控制权}账本记录曹魏与晋王国在咸熙二年八月（265年）发生“司马昭去世，司马炎继承晋王位及魏廷实际军政控制权”。条目把背景概括为司马昭已完成晋王封授和灭蜀，却在受禅前死亡，司马氏必须避免权力交接给曹魏残余以可乘之机，并指出司马炎迅速承接晋王国与魏廷资源，为同年正式取代曹魏创造连续性；家族继权不等于自动获得礼仪合法性。这个节点进入区域社会时，要经由文书、道路、军队、市场、寺院和家庭等不同中介；因此，政治行动的宣布、地方接收和日常生活的改变不能被视作同一时刻完成。

本条依据《三国志·陈留王纪》；《晋书·文帝纪》；《晋书·武帝纪》，证据提示为“卒年与继承可靠；遗命、司马炎接掌禁军和魏帝朝臣的协商空间有争议”。这些材料可支持年序、官方决定、战事或特定地点的线索，却难以直接统计所有人的财富、迁徙、信仰和动机。更稳妥的读法是区分诏令与实施、军报与人口、行政称谓与自我认同，并让地方材料的缺环限制解释的范围。

由此可见，该事件不是孤立的年表点，而是资源、信息与权力关系重新配置的一环。不同地区的官员、社群和家庭可能以不同方式承担征发、保护、迁移或宗教活动；现存来源只能照亮其中一部分，不能替沉默者制造统一的经验。

\eventanchor{TK-0265-WEI-ABDICATION}{咸熙二年十二月（265年）·曹奂禅位于司马炎，西晋建立，曹魏皇帝体系终结}账本记录曹魏与西晋在咸熙二年十二月（265年）发生“曹奂禅位于司马炎，西晋建立，曹魏皇帝体系终结”。条目把背景概括为司马氏已控制魏的军队、诏令和封国机构，灭蜀后又消除最强的内部军事竞争者，并指出西晋继承魏的北方行政与军事资源，并将统一目标转向孙吴；禅让文书应视为既成权力转移的礼仪化表达。这个节点进入区域社会时，要经由文书、道路、军队、市场、寺院和家庭等不同中介；因此，政治行动的宣布、地方接收和日常生活的改变不能被视作同一时刻完成。

本条依据《三国志·陈留王纪》；《晋书·武帝纪》；《资治通鉴》晋纪，证据提示为“禅代年月可靠；诏册撰作、群臣参与和曹奂的选择空间有争议”。这些材料可支持年序、官方决定、战事或特定地点的线索，却难以直接统计所有人的财富、迁徙、信仰和动机。更稳妥的读法是区分诏令与实施、军报与人口、行政称谓与自我认同，并让地方材料的缺环限制解释的范围。

由此可见，该事件不是孤立的年表点，而是资源、信息与权力关系重新配置的一环。不同地区的官员、社群和家庭可能以不同方式承担征发、保护、迁移或宗教活动；现存来源只能照亮其中一部分，不能替沉默者制造统一的经验。

\eventanchor{TK-0268-TAISHI-CODE}{泰始四年（268年）·西晋完成并颁行《泰始律》，以十八篇新律整理魏晋以来的刑法框架}账本记录西晋在泰始四年（268年）发生“西晋完成并颁行《泰始律》，以十八篇新律整理魏晋以来的刑法框架”。条目把背景概括为西晋取代魏后需要将军政接管转化为可持续的官僚规范，并处理魏末以来律令繁杂的问题，并指出新律为统一前后的司法提供制度语言，但法典条文不等同于各州郡审判实践；后世法典只能作为传承线索而非原文替代。这个节点进入区域社会时，要经由文书、道路、军队、市场、寺院和家庭等不同中介；因此，政治行动的宣布、地方接收和日常生活的改变不能被视作同一时刻完成。

本条依据《晋书·刑法志》；《晋书·武帝纪》；《唐律疏议》所见后世法律史线索，证据提示为“成律年代与篇数大体可靠；实际颁行范围、与旧令关系和地方执行程度有争议”。这些材料可支持年序、官方决定、战事或特定地点的线索，却难以直接统计所有人的财富、迁徙、信仰和动机。更稳妥的读法是区分诏令与实施、军报与人口、行政称谓与自我认同，并让地方材料的缺环限制解释的范围。

由此可见，该事件不是孤立的年表点，而是资源、信息与权力关系重新配置的一环。不同地区的官员、社群和家庭可能以不同方式承担征发、保护、迁移或宗教活动；现存来源只能照亮其中一部分，不能替沉默者制造统一的经验。

\eventanchor{TK-0269-YANG-HU-JINGZHOU}{泰始五年（269年）·羊祜都督荆州诸军事，在襄阳一线经营屯戍、招纳与对吴关系}账本记录西晋与孙吴在泰始五年（269年）发生“羊祜都督荆州诸军事，在襄阳一线经营屯戍、招纳与对吴关系”。条目把背景概括为西晋灭蜀后与孙吴直接相邻，长江中游既需防吴军北上，也需积累渡江和安置降附者的条件，并指出襄阳成为晋的前进基地，羊祜与陆抗之间的边境互动影响后来伐吴的路线与政治宣传；招纳数字须审慎对待。这个节点进入区域社会时，要经由文书、道路、军队、市场、寺院和家庭等不同中介；因此，政治行动的宣布、地方接收和日常生活的改变不能被视作同一时刻完成。

本条依据《晋书·羊祜传》；《三国志·陆抗传》；《资治通鉴》晋纪，证据提示为“任命与荆州经营可靠；招降规模、怀柔政策效果和羊祜个人作用有争议”。这些材料可支持年序、官方决定、战事或特定地点的线索，却难以直接统计所有人的财富、迁徙、信仰和动机。更稳妥的读法是区分诏令与实施、军报与人口、行政称谓与自我认同，并让地方材料的缺环限制解释的范围。

由此可见，该事件不是孤立的年表点，而是资源、信息与权力关系重新配置的一环。不同地区的官员、社群和家庭可能以不同方式承担征发、保护、迁移或宗教活动；现存来源只能照亮其中一部分，不能替沉默者制造统一的经验。

\eventanchor{TK-0272-XILING}{凤凰元年（272年）·西陵督步阐降晋，陆抗围西陵并击退晋援，孙吴收复峡口要地}账本记录孙吴与西晋在凤凰元年（272年）发生“西陵督步阐降晋，陆抗围西陵并击退晋援，孙吴收复峡口要地”。条目把背景概括为步阐与孙吴朝廷关系恶化，西晋希望利用内应取得长江上游突破口，陆抗则须守住武昌以西的水陆枢纽，并指出孙吴暂时阻止晋军楔入三峡—江汉通道，但西陵危机暴露吴内部将领与中央的脆弱联系，晋继续修订伐吴方案。这个节点进入区域社会时，要经由文书、道路、军队、市场、寺院和家庭等不同中介；因此，政治行动的宣布、地方接收和日常生活的改变不能被视作同一时刻完成。

本条依据《三国志·陆抗传》；《晋书·羊祜传》；《晋书·杜预传》；《资治通鉴》晋纪，证据提示为“步阐降晋、陆抗取胜和西陵收复可靠；晋援军协同、围城经过和双方损失有争议”。这些材料可支持年序、官方决定、战事或特定地点的线索，却难以直接统计所有人的财富、迁徙、信仰和动机。更稳妥的读法是区分诏令与实施、军报与人口、行政称谓与自我认同，并让地方材料的缺环限制解释的范围。

由此可见，该事件不是孤立的年表点，而是资源、信息与权力关系重新配置的一环。不同地区的官员、社群和家庭可能以不同方式承担征发、保护、迁移或宗教活动；现存来源只能照亮其中一部分，不能替沉默者制造统一的经验。

\eventanchor{TK-0274-LU-KANG-DEATH}{凤凰三年（274年）·陆抗去世，孙吴失去镇守荆州西部、熟悉晋吴边境的主要统帅}账本记录孙吴在凤凰三年（274年）发生“陆抗去世，孙吴失去镇守荆州西部、熟悉晋吴边境的主要统帅”。条目把背景概括为陆抗在西陵之役后维系荆州防线，其军政经验对孙吴抵御晋的上游压力至关重要，并指出孙吴难以用一名统帅完整替代其边境网络，西部防线的协调能力下降；遗表中的战略主张须注意文本传承和追述成分。这个节点进入区域社会时，要经由文书、道路、军队、市场、寺院和家庭等不同中介；因此，政治行动的宣布、地方接收和日常生活的改变不能被视作同一时刻完成。

本条依据《三国志·陆抗传》；《建康实录》；《晋书·羊祜传》，证据提示为“卒年和官职可靠；病因、其遗表传本及死后荆州部署变化有争议”。这些材料可支持年序、官方决定、战事或特定地点的线索，却难以直接统计所有人的财富、迁徙、信仰和动机。更稳妥的读法是区分诏令与实施、军报与人口、行政称谓与自我认同，并让地方材料的缺环限制解释的范围。

由此可见，该事件不是孤立的年表点，而是资源、信息与权力关系重新配置的一环。不同地区的官员、社群和家庭可能以不同方式承担征发、保护、迁移或宗教活动；现存来源只能照亮其中一部分，不能替沉默者制造统一的经验。

\eventanchor{TK-0278-YANG-HU-DEATH}{咸宁四年（278年）·羊祜去世前上表建议伐吴，杜预随后接掌荆州军事部署}账本记录西晋在咸宁四年（278年）发生“羊祜去世前上表建议伐吴，杜预随后接掌荆州军事部署”。条目把背景概括为西晋已在襄阳、益州和淮南积累进攻条件，羊祜担心长期对峙会给孙吴修复边防的机会，并指出杜预继承荆州方面的准备，朝廷由长期边防经营转向具体的多路伐吴计划；不能把后来的统一完全归因于一篇遗表。这个节点进入区域社会时，要经由文书、道路、军队、市场、寺院和家庭等不同中介；因此，政治行动的宣布、地方接收和日常生活的改变不能被视作同一时刻完成。

本条依据《晋书·羊祜传》；《晋书·杜预传》；《资治通鉴》晋纪，证据提示为“羊祜卒、表请伐吴和杜预继任大体可靠；表文传本、羊祜战略影响和朝廷决策先后有争议”。这些材料可支持年序、官方决定、战事或特定地点的线索，却难以直接统计所有人的财富、迁徙、信仰和动机。更稳妥的读法是区分诏令与实施、军报与人口、行政称谓与自我认同，并让地方材料的缺环限制解释的范围。

由此可见，该事件不是孤立的年表点，而是资源、信息与权力关系重新配置的一环。不同地区的官员、社群和家庭可能以不同方式承担征发、保护、迁移或宗教活动；现存来源只能照亮其中一部分，不能替沉默者制造统一的经验。

\eventanchor{TK-0279-JIN-INVADES-WU}{咸宁五年十一月（279年）·西晋以杜预、王濬、王浑等分路伐吴，益州水军沿江东下，荆淮军并进}账本记录西晋与孙吴在咸宁五年十一月（279年）发生“西晋以杜预、王濬、王浑等分路伐吴，益州水军沿江东下，荆淮军并进”。条目把背景概括为晋在灭蜀、整顿法制和长期经营荆淮后，判断孙吴朝廷内耗、边将离心且难以同时应对水陆攻势，并指出战争突破三国长期均势，西晋舰队和沿江攻城将决定建业能否获得外围防线和地方将领的支持。这个节点进入区域社会时，要经由文书、道路、军队、市场、寺院和家庭等不同中介；因此，政治行动的宣布、地方接收和日常生活的改变不能被视作同一时刻完成。

本条依据《晋书·武帝纪》；《晋书·杜预传》；《晋书·王濬传》；《三国志·孙皓传》；《资治通鉴》晋纪，证据提示为“出师时间、主要统帅和多路部署可靠；兵力数字、路线细节、各军协同与王濬舰队规模有争议”。这些材料可支持年序、官方决定、战事或特定地点的线索，却难以直接统计所有人的财富、迁徙、信仰和动机。更稳妥的读法是区分诏令与实施、军报与人口、行政称谓与自我认同，并让地方材料的缺环限制解释的范围。

由此可见，该事件不是孤立的年表点，而是资源、信息与权力关系重新配置的一环。不同地区的官员、社群和家庭可能以不同方式承担征发、保护、迁移或宗教活动；现存来源只能照亮其中一部分，不能替沉默者制造统一的经验。

\eventanchor{TK-0280-SUN-HAO-SURRENDER}{太康元年三月（280年）·王濬舰队抵建业，孙皓出降，孙吴灭亡，西晋完成全国统一}账本记录西晋与孙吴在太康元年三月（280年）发生“王濬舰队抵建业，孙皓出降，孙吴灭亡，西晋完成全国统一”。条目把背景概括为晋军沿江连续突破，吴的地方将领和中央难以重建防线，建业在军事与政治上同时孤立，并指出自220年以来的三国并立结束，西晋接收吴的州郡与官民；统一不等于地方秩序、赋役和江南社会立即无缝整合。这个节点进入区域社会时，要经由文书、道路、军队、市场、寺院和家庭等不同中介；因此，政治行动的宣布、地方接收和日常生活的改变不能被视作同一时刻完成。

本条依据《晋书·武帝纪》；《晋书·王濬传》；《三国志·孙皓传》；《建康实录》，证据提示为“建业降附、孙吴灭亡和统一结果可靠；王濬与王浑的功绩争议、降仪细节、战后安置和统一的实际覆盖范围有争议”。这些材料可支持年序、官方决定、战事或特定地点的线索，却难以直接统计所有人的财富、迁徙、信仰和动机。更稳妥的读法是区分诏令与实施、军报与人口、行政称谓与自我认同，并让地方材料的缺环限制解释的范围。

由此可见，该事件不是孤立的年表点，而是资源、信息与权力关系重新配置的一环。不同地区的官员、社群和家庭可能以不同方式承担征发、保护、迁移或宗教活动；现存来源只能照亮其中一部分，不能替沉默者制造统一的经验。

\eventanchor{TK-LEDGER-001-1}{黄初元年（220年）·“曹丕定都洛阳，以魏的尚书、禁军和北方仓储网络承接汉末朝廷”的前期动员与资源准备}账本记录曹魏在黄初元年（220年）发生““曹丕定都洛阳，以魏的尚书、禁军和北方仓储网络承接汉末朝廷”的前期动员与资源准备”。条目把背景概括为当时的权力、资源与信息约束推动相关过程展开，并指出这一过程为“曹丕定都洛阳，以魏的尚书、禁军和北方仓储网络承接汉末朝廷”提供可观察的条件。这个节点进入区域社会时，要经由文书、道路、军队、市场、寺院和家庭等不同中介；因此，政治行动的宣布、地方接收和日常生活的改变不能被视作同一时刻完成。

本条依据《三国志·文帝纪》；《三国志·明帝纪》；《资治通鉴》魏纪，证据提示为“核心事项已有既有账本来源；本分解节点的参与者、执行范围和时间先后仍须逐案核验”。这些材料可支持年序、官方决定、战事或特定地点的线索，却难以直接统计所有人的财富、迁徙、信仰和动机。更稳妥的读法是区分诏令与实施、军报与人口、行政称谓与自我认同，并让地方材料的缺环限制解释的范围。

由此可见，该事件不是孤立的年表点，而是资源、信息与权力关系重新配置的一环。不同地区的官员、社群和家庭可能以不同方式承担征发、保护、迁移或宗教活动；现存来源只能照亮其中一部分，不能替沉默者制造统一的经验。

\eventanchor{TK-LEDGER-001-2}{黄初元年（220年）·“曹丕定都洛阳，以魏的尚书、禁军和北方仓储网络承接汉末朝廷”的命令、联盟或地方响应}账本记录曹魏在黄初元年（220年）发生““曹丕定都洛阳，以魏的尚书、禁军和北方仓储网络承接汉末朝廷”的命令、联盟或地方响应”。条目把背景概括为当时的权力、资源与信息约束推动相关过程展开，并指出这一过程为“曹丕定都洛阳，以魏的尚书、禁军和北方仓储网络承接汉末朝廷”提供可观察的条件。这个节点进入区域社会时，要经由文书、道路、军队、市场、寺院和家庭等不同中介；因此，政治行动的宣布、地方接收和日常生活的改变不能被视作同一时刻完成。

本条依据《三国志·文帝纪》；《三国志·明帝纪》；《资治通鉴》魏纪，证据提示为“核心事项已有既有账本来源；本分解节点的参与者、执行范围和时间先后仍须逐案核验”。这些材料可支持年序、官方决定、战事或特定地点的线索，却难以直接统计所有人的财富、迁徙、信仰和动机。更稳妥的读法是区分诏令与实施、军报与人口、行政称谓与自我认同，并让地方材料的缺环限制解释的范围。

由此可见，该事件不是孤立的年表点，而是资源、信息与权力关系重新配置的一环。不同地区的官员、社群和家庭可能以不同方式承担征发、保护、迁移或宗教活动；现存来源只能照亮其中一部分，不能替沉默者制造统一的经验。

\subsection{边疆、战争与移动}

\eventanchor{TK-LEDGER-001-3}{黄初元年（220年）·曹丕定都洛阳，以魏的尚书、禁军和北方仓储网络承接汉末朝廷}账本记录曹魏在黄初元年（220年）发生“曹丕定都洛阳，以魏的尚书、禁军和北方仓储网络承接汉末朝廷”。条目把背景概括为“曹丕定都洛阳，以魏的尚书、禁军和北方仓储网络承接汉末朝廷”发生的政治与区域条件共同作用，并指出该节点改变后续资源、联盟或制度处置的条件。这个节点进入区域社会时，要经由文书、道路、军队、市场、寺院和家庭等不同中介；因此，政治行动的宣布、地方接收和日常生活的改变不能被视作同一时刻完成。

本条依据《三国志·文帝纪》；《三国志·明帝纪》；《资治通鉴》魏纪，证据提示为“核心事项已有既有账本来源；本分解节点的参与者、执行范围和时间先后仍须逐案核验”。这些材料可支持年序、官方决定、战事或特定地点的线索，却难以直接统计所有人的财富、迁徙、信仰和动机。更稳妥的读法是区分诏令与实施、军报与人口、行政称谓与自我认同，并让地方材料的缺环限制解释的范围。

由此可见，该事件不是孤立的年表点，而是资源、信息与权力关系重新配置的一环。不同地区的官员、社群和家庭可能以不同方式承担征发、保护、迁移或宗教活动；现存来源只能照亮其中一部分，不能替沉默者制造统一的经验。

\eventanchor{TK-LEDGER-001-4}{黄初元年（220年）·“曹丕定都洛阳，以魏的尚书、禁军和北方仓储网络承接汉末朝廷”的执行中的空间与人员处置}账本记录曹魏在黄初元年（220年）发生““曹丕定都洛阳，以魏的尚书、禁军和北方仓储网络承接汉末朝廷”的执行中的空间与人员处置”。条目把背景概括为“曹丕定都洛阳，以魏的尚书、禁军和北方仓储网络承接汉末朝廷”发生的政治与区域条件共同作用，并指出该节点改变后续资源、联盟或制度处置的条件。这个节点进入区域社会时，要经由文书、道路、军队、市场、寺院和家庭等不同中介；因此，政治行动的宣布、地方接收和日常生活的改变不能被视作同一时刻完成。

本条依据《三国志·文帝纪》；《三国志·明帝纪》；《资治通鉴》魏纪，证据提示为“核心事项已有既有账本来源；本分解节点的参与者、执行范围和时间先后仍须逐案核验”。这些材料可支持年序、官方决定、战事或特定地点的线索，却难以直接统计所有人的财富、迁徙、信仰和动机。更稳妥的读法是区分诏令与实施、军报与人口、行政称谓与自我认同，并让地方材料的缺环限制解释的范围。

由此可见，该事件不是孤立的年表点，而是资源、信息与权力关系重新配置的一环。不同地区的官员、社群和家庭可能以不同方式承担征发、保护、迁移或宗教活动；现存来源只能照亮其中一部分，不能替沉默者制造统一的经验。

\eventanchor{TK-LEDGER-001-5}{黄初元年（220年）·“曹丕定都洛阳，以魏的尚书、禁军和北方仓储网络承接汉末朝廷”的后续制度或军政调整}账本记录曹魏在黄初元年（220年）发生““曹丕定都洛阳，以魏的尚书、禁军和北方仓储网络承接汉末朝廷”的后续制度或军政调整”。条目把背景概括为核心节点发生后仍须处理其遗留问题，并指出后续影响须在不同地区和材料类型中分别核验。这个节点进入区域社会时，要经由文书、道路、军队、市场、寺院和家庭等不同中介；因此，政治行动的宣布、地方接收和日常生活的改变不能被视作同一时刻完成。

本条依据《三国志·文帝纪》；《三国志·明帝纪》；《资治通鉴》魏纪，证据提示为“核心事项已有既有账本来源；本分解节点的参与者、执行范围和时间先后仍须逐案核验”。这些材料可支持年序、官方决定、战事或特定地点的线索，却难以直接统计所有人的财富、迁徙、信仰和动机。更稳妥的读法是区分诏令与实施、军报与人口、行政称谓与自我认同，并让地方材料的缺环限制解释的范围。

由此可见，该事件不是孤立的年表点，而是资源、信息与权力关系重新配置的一环。不同地区的官员、社群和家庭可能以不同方式承担征发、保护、迁移或宗教活动；现存来源只能照亮其中一部分，不能替沉默者制造统一的经验。

\eventanchor{TK-LEDGER-001-6}{黄初元年（220年）·“曹丕定都洛阳，以魏的尚书、禁军和北方仓储网络承接汉末朝廷”的异源材料与影响范围的复核}账本记录曹魏在黄初元年（220年）发生““曹丕定都洛阳，以魏的尚书、禁军和北方仓储网络承接汉末朝廷”的异源材料与影响范围的复核”。条目把背景概括为核心节点发生后仍须处理其遗留问题，并指出后续影响须在不同地区和材料类型中分别核验。这个节点进入区域社会时，要经由文书、道路、军队、市场、寺院和家庭等不同中介；因此，政治行动的宣布、地方接收和日常生活的改变不能被视作同一时刻完成。

本条依据《三国志·文帝纪》；《三国志·明帝纪》；《资治通鉴》魏纪，证据提示为“核心事项已有既有账本来源；本分解节点的参与者、执行范围和时间先后仍须逐案核验”。这些材料可支持年序、官方决定、战事或特定地点的线索，却难以直接统计所有人的财富、迁徙、信仰和动机。更稳妥的读法是区分诏令与实施、军报与人口、行政称谓与自我认同，并让地方材料的缺环限制解释的范围。

由此可见，该事件不是孤立的年表点，而是资源、信息与权力关系重新配置的一环。不同地区的官员、社群和家庭可能以不同方式承担征发、保护、迁移或宗教活动；现存来源只能照亮其中一部分，不能替沉默者制造统一的经验。

\eventanchor{TK-LEDGER-002-1}{章武元年四月（221年）·“刘备在成都即皇帝位，建元章武，以汉室继承者自居”的前期动员与资源准备}账本记录蜀汉在章武元年四月（221年）发生““刘备在成都即皇帝位，建元章武，以汉室继承者自居”的前期动员与资源准备”。条目把背景概括为当时的权力、资源与信息约束推动相关过程展开，并指出这一过程为“刘备在成都即皇帝位，建元章武，以汉室继承者自居”提供可观察的条件。这个节点进入区域社会时，要经由文书、道路、军队、市场、寺院和家庭等不同中介；因此，政治行动的宣布、地方接收和日常生活的改变不能被视作同一时刻完成。

本条依据《三国志·先主传》；《三国志·后主传》；《华阳国志·刘先主志》，证据提示为“核心事项已有既有账本来源；本分解节点的参与者、执行范围和时间先后仍须逐案核验”。这些材料可支持年序、官方决定、战事或特定地点的线索，却难以直接统计所有人的财富、迁徙、信仰和动机。更稳妥的读法是区分诏令与实施、军报与人口、行政称谓与自我认同，并让地方材料的缺环限制解释的范围。

由此可见，该事件不是孤立的年表点，而是资源、信息与权力关系重新配置的一环。不同地区的官员、社群和家庭可能以不同方式承担征发、保护、迁移或宗教活动；现存来源只能照亮其中一部分，不能替沉默者制造统一的经验。

\eventanchor{TK-LEDGER-002-2}{章武元年四月（221年）·“刘备在成都即皇帝位，建元章武，以汉室继承者自居”的命令、联盟或地方响应}账本记录蜀汉在章武元年四月（221年）发生““刘备在成都即皇帝位，建元章武，以汉室继承者自居”的命令、联盟或地方响应”。条目把背景概括为当时的权力、资源与信息约束推动相关过程展开，并指出这一过程为“刘备在成都即皇帝位，建元章武，以汉室继承者自居”提供可观察的条件。这个节点进入区域社会时，要经由文书、道路、军队、市场、寺院和家庭等不同中介；因此，政治行动的宣布、地方接收和日常生活的改变不能被视作同一时刻完成。

本条依据《三国志·先主传》；《三国志·后主传》；《华阳国志·刘先主志》，证据提示为“核心事项已有既有账本来源；本分解节点的参与者、执行范围和时间先后仍须逐案核验”。这些材料可支持年序、官方决定、战事或特定地点的线索，却难以直接统计所有人的财富、迁徙、信仰和动机。更稳妥的读法是区分诏令与实施、军报与人口、行政称谓与自我认同，并让地方材料的缺环限制解释的范围。

由此可见，该事件不是孤立的年表点，而是资源、信息与权力关系重新配置的一环。不同地区的官员、社群和家庭可能以不同方式承担征发、保护、迁移或宗教活动；现存来源只能照亮其中一部分，不能替沉默者制造统一的经验。

\eventanchor{TK-LEDGER-002-3}{章武元年四月（221年）·刘备在成都即皇帝位，建元章武，以汉室继承者自居}账本记录蜀汉在章武元年四月（221年）发生“刘备在成都即皇帝位，建元章武，以汉室继承者自居”。条目把背景概括为“刘备在成都即皇帝位，建元章武，以汉室继承者自居”发生的政治与区域条件共同作用，并指出该节点改变后续资源、联盟或制度处置的条件。这个节点进入区域社会时，要经由文书、道路、军队、市场、寺院和家庭等不同中介；因此，政治行动的宣布、地方接收和日常生活的改变不能被视作同一时刻完成。

本条依据《三国志·先主传》；《三国志·后主传》；《华阳国志·刘先主志》，证据提示为“核心事项已有既有账本来源；本分解节点的参与者、执行范围和时间先后仍须逐案核验”。这些材料可支持年序、官方决定、战事或特定地点的线索，却难以直接统计所有人的财富、迁徙、信仰和动机。更稳妥的读法是区分诏令与实施、军报与人口、行政称谓与自我认同，并让地方材料的缺环限制解释的范围。

由此可见，该事件不是孤立的年表点，而是资源、信息与权力关系重新配置的一环。不同地区的官员、社群和家庭可能以不同方式承担征发、保护、迁移或宗教活动；现存来源只能照亮其中一部分，不能替沉默者制造统一的经验。

\eventanchor{TK-LEDGER-002-4}{章武元年四月（221年）·“刘备在成都即皇帝位，建元章武，以汉室继承者自居”的执行中的空间与人员处置}账本记录蜀汉在章武元年四月（221年）发生““刘备在成都即皇帝位，建元章武，以汉室继承者自居”的执行中的空间与人员处置”。条目把背景概括为“刘备在成都即皇帝位，建元章武，以汉室继承者自居”发生的政治与区域条件共同作用，并指出该节点改变后续资源、联盟或制度处置的条件。这个节点进入区域社会时，要经由文书、道路、军队、市场、寺院和家庭等不同中介；因此，政治行动的宣布、地方接收和日常生活的改变不能被视作同一时刻完成。

本条依据《三国志·先主传》；《三国志·后主传》；《华阳国志·刘先主志》，证据提示为“核心事项已有既有账本来源；本分解节点的参与者、执行范围和时间先后仍须逐案核验”。这些材料可支持年序、官方决定、战事或特定地点的线索，却难以直接统计所有人的财富、迁徙、信仰和动机。更稳妥的读法是区分诏令与实施、军报与人口、行政称谓与自我认同，并让地方材料的缺环限制解释的范围。

由此可见，该事件不是孤立的年表点，而是资源、信息与权力关系重新配置的一环。不同地区的官员、社群和家庭可能以不同方式承担征发、保护、迁移或宗教活动；现存来源只能照亮其中一部分，不能替沉默者制造统一的经验。

\eventanchor{TK-LEDGER-002-5}{章武元年四月（221年）·“刘备在成都即皇帝位，建元章武，以汉室继承者自居”的后续制度或军政调整}账本记录蜀汉在章武元年四月（221年）发生““刘备在成都即皇帝位，建元章武，以汉室继承者自居”的后续制度或军政调整”。条目把背景概括为核心节点发生后仍须处理其遗留问题，并指出后续影响须在不同地区和材料类型中分别核验。这个节点进入区域社会时，要经由文书、道路、军队、市场、寺院和家庭等不同中介；因此，政治行动的宣布、地方接收和日常生活的改变不能被视作同一时刻完成。

本条依据《三国志·先主传》；《三国志·后主传》；《华阳国志·刘先主志》，证据提示为“核心事项已有既有账本来源；本分解节点的参与者、执行范围和时间先后仍须逐案核验”。这些材料可支持年序、官方决定、战事或特定地点的线索，却难以直接统计所有人的财富、迁徙、信仰和动机。更稳妥的读法是区分诏令与实施、军报与人口、行政称谓与自我认同，并让地方材料的缺环限制解释的范围。

由此可见，该事件不是孤立的年表点，而是资源、信息与权力关系重新配置的一环。不同地区的官员、社群和家庭可能以不同方式承担征发、保护、迁移或宗教活动；现存来源只能照亮其中一部分，不能替沉默者制造统一的经验。

\eventanchor{TK-LEDGER-002-6}{章武元年四月（221年）·“刘备在成都即皇帝位，建元章武，以汉室继承者自居”的异源材料与影响范围的复核}账本记录蜀汉在章武元年四月（221年）发生““刘备在成都即皇帝位，建元章武，以汉室继承者自居”的异源材料与影响范围的复核”。条目把背景概括为核心节点发生后仍须处理其遗留问题，并指出后续影响须在不同地区和材料类型中分别核验。这个节点进入区域社会时，要经由文书、道路、军队、市场、寺院和家庭等不同中介；因此，政治行动的宣布、地方接收和日常生活的改变不能被视作同一时刻完成。

本条依据《三国志·先主传》；《三国志·后主传》；《华阳国志·刘先主志》，证据提示为“核心事项已有既有账本来源；本分解节点的参与者、执行范围和时间先后仍须逐案核验”。这些材料可支持年序、官方决定、战事或特定地点的线索，却难以直接统计所有人的财富、迁徙、信仰和动机。更稳妥的读法是区分诏令与实施、军报与人口、行政称谓与自我认同，并让地方材料的缺环限制解释的范围。

由此可见，该事件不是孤立的年表点，而是资源、信息与权力关系重新配置的一环。不同地区的官员、社群和家庭可能以不同方式承担征发、保护、迁移或宗教活动；现存来源只能照亮其中一部分，不能替沉默者制造统一的经验。

\eventanchor{TK-LEDGER-003-1}{章武元年（221年）·“刘备集结益州军东下，准备以夺回荆州和为关羽报仇为由进攻孙吴”的前期动员与资源准备}账本记录蜀汉与孙吴在章武元年（221年）发生““刘备集结益州军东下，准备以夺回荆州和为关羽报仇为由进攻孙吴”的前期动员与资源准备”。条目把背景概括为当时的权力、资源与信息约束推动相关过程展开，并指出这一过程为“刘备集结益州军东下，准备以夺回荆州和为关羽报仇为由进攻孙吴”提供可观察的条件。这个节点进入区域社会时，要经由文书、道路、军队、市场、寺院和家庭等不同中介；因此，政治行动的宣布、地方接收和日常生活的改变不能被视作同一时刻完成。

本条依据《三国志·先主传》；《三国志·张飞传》；《三国志·陆逊传》，证据提示为“核心事项已有既有账本来源；本分解节点的参与者、执行范围和时间先后仍须逐案核验”。这些材料可支持年序、官方决定、战事或特定地点的线索，却难以直接统计所有人的财富、迁徙、信仰和动机。更稳妥的读法是区分诏令与实施、军报与人口、行政称谓与自我认同，并让地方材料的缺环限制解释的范围。

由此可见，该事件不是孤立的年表点，而是资源、信息与权力关系重新配置的一环。不同地区的官员、社群和家庭可能以不同方式承担征发、保护、迁移或宗教活动；现存来源只能照亮其中一部分，不能替沉默者制造统一的经验。

\eventanchor{TK-LEDGER-003-2}{章武元年（221年）·“刘备集结益州军东下，准备以夺回荆州和为关羽报仇为由进攻孙吴”的命令、联盟或地方响应}账本记录蜀汉与孙吴在章武元年（221年）发生““刘备集结益州军东下，准备以夺回荆州和为关羽报仇为由进攻孙吴”的命令、联盟或地方响应”。条目把背景概括为当时的权力、资源与信息约束推动相关过程展开，并指出这一过程为“刘备集结益州军东下，准备以夺回荆州和为关羽报仇为由进攻孙吴”提供可观察的条件。这个节点进入区域社会时，要经由文书、道路、军队、市场、寺院和家庭等不同中介；因此，政治行动的宣布、地方接收和日常生活的改变不能被视作同一时刻完成。

本条依据《三国志·先主传》；《三国志·张飞传》；《三国志·陆逊传》，证据提示为“核心事项已有既有账本来源；本分解节点的参与者、执行范围和时间先后仍须逐案核验”。这些材料可支持年序、官方决定、战事或特定地点的线索，却难以直接统计所有人的财富、迁徙、信仰和动机。更稳妥的读法是区分诏令与实施、军报与人口、行政称谓与自我认同，并让地方材料的缺环限制解释的范围。

由此可见，该事件不是孤立的年表点，而是资源、信息与权力关系重新配置的一环。不同地区的官员、社群和家庭可能以不同方式承担征发、保护、迁移或宗教活动；现存来源只能照亮其中一部分，不能替沉默者制造统一的经验。

\eventanchor{TK-LEDGER-003-3}{章武元年（221年）·刘备集结益州军东下，准备以夺回荆州和为关羽报仇为由进攻孙吴}账本记录蜀汉与孙吴在章武元年（221年）发生“刘备集结益州军东下，准备以夺回荆州和为关羽报仇为由进攻孙吴”。条目把背景概括为“刘备集结益州军东下，准备以夺回荆州和为关羽报仇为由进攻孙吴”发生的政治与区域条件共同作用，并指出该节点改变后续资源、联盟或制度处置的条件。这个节点进入区域社会时，要经由文书、道路、军队、市场、寺院和家庭等不同中介；因此，政治行动的宣布、地方接收和日常生活的改变不能被视作同一时刻完成。

本条依据《三国志·先主传》；《三国志·张飞传》；《三国志·陆逊传》，证据提示为“核心事项已有既有账本来源；本分解节点的参与者、执行范围和时间先后仍须逐案核验”。这些材料可支持年序、官方决定、战事或特定地点的线索，却难以直接统计所有人的财富、迁徙、信仰和动机。更稳妥的读法是区分诏令与实施、军报与人口、行政称谓与自我认同，并让地方材料的缺环限制解释的范围。

由此可见，该事件不是孤立的年表点，而是资源、信息与权力关系重新配置的一环。不同地区的官员、社群和家庭可能以不同方式承担征发、保护、迁移或宗教活动；现存来源只能照亮其中一部分，不能替沉默者制造统一的经验。

\eventanchor{TK-LEDGER-003-4}{章武元年（221年）·“刘备集结益州军东下，准备以夺回荆州和为关羽报仇为由进攻孙吴”的执行中的空间与人员处置}账本记录蜀汉与孙吴在章武元年（221年）发生““刘备集结益州军东下，准备以夺回荆州和为关羽报仇为由进攻孙吴”的执行中的空间与人员处置”。条目把背景概括为“刘备集结益州军东下，准备以夺回荆州和为关羽报仇为由进攻孙吴”发生的政治与区域条件共同作用，并指出该节点改变后续资源、联盟或制度处置的条件。这个节点进入区域社会时，要经由文书、道路、军队、市场、寺院和家庭等不同中介；因此，政治行动的宣布、地方接收和日常生活的改变不能被视作同一时刻完成。

本条依据《三国志·先主传》；《三国志·张飞传》；《三国志·陆逊传》，证据提示为“核心事项已有既有账本来源；本分解节点的参与者、执行范围和时间先后仍须逐案核验”。这些材料可支持年序、官方决定、战事或特定地点的线索，却难以直接统计所有人的财富、迁徙、信仰和动机。更稳妥的读法是区分诏令与实施、军报与人口、行政称谓与自我认同，并让地方材料的缺环限制解释的范围。

由此可见，该事件不是孤立的年表点，而是资源、信息与权力关系重新配置的一环。不同地区的官员、社群和家庭可能以不同方式承担征发、保护、迁移或宗教活动；现存来源只能照亮其中一部分，不能替沉默者制造统一的经验。

\eventanchor{TK-LEDGER-003-5}{章武元年（221年）·“刘备集结益州军东下，准备以夺回荆州和为关羽报仇为由进攻孙吴”的后续制度或军政调整}账本记录蜀汉与孙吴在章武元年（221年）发生““刘备集结益州军东下，准备以夺回荆州和为关羽报仇为由进攻孙吴”的后续制度或军政调整”。条目把背景概括为核心节点发生后仍须处理其遗留问题，并指出后续影响须在不同地区和材料类型中分别核验。这个节点进入区域社会时，要经由文书、道路、军队、市场、寺院和家庭等不同中介；因此，政治行动的宣布、地方接收和日常生活的改变不能被视作同一时刻完成。

本条依据《三国志·先主传》；《三国志·张飞传》；《三国志·陆逊传》，证据提示为“核心事项已有既有账本来源；本分解节点的参与者、执行范围和时间先后仍须逐案核验”。这些材料可支持年序、官方决定、战事或特定地点的线索，却难以直接统计所有人的财富、迁徙、信仰和动机。更稳妥的读法是区分诏令与实施、军报与人口、行政称谓与自我认同，并让地方材料的缺环限制解释的范围。

由此可见，该事件不是孤立的年表点，而是资源、信息与权力关系重新配置的一环。不同地区的官员、社群和家庭可能以不同方式承担征发、保护、迁移或宗教活动；现存来源只能照亮其中一部分，不能替沉默者制造统一的经验。

\eventanchor{TK-LEDGER-003-6}{章武元年（221年）·“刘备集结益州军东下，准备以夺回荆州和为关羽报仇为由进攻孙吴”的异源材料与影响范围的复核}账本记录蜀汉与孙吴在章武元年（221年）发生““刘备集结益州军东下，准备以夺回荆州和为关羽报仇为由进攻孙吴”的异源材料与影响范围的复核”。条目把背景概括为核心节点发生后仍须处理其遗留问题，并指出后续影响须在不同地区和材料类型中分别核验。这个节点进入区域社会时，要经由文书、道路、军队、市场、寺院和家庭等不同中介；因此，政治行动的宣布、地方接收和日常生活的改变不能被视作同一时刻完成。

本条依据《三国志·先主传》；《三国志·张飞传》；《三国志·陆逊传》，证据提示为“核心事项已有既有账本来源；本分解节点的参与者、执行范围和时间先后仍须逐案核验”。这些材料可支持年序、官方决定、战事或特定地点的线索，却难以直接统计所有人的财富、迁徙、信仰和动机。更稳妥的读法是区分诏令与实施、军报与人口、行政称谓与自我认同，并让地方材料的缺环限制解释的范围。

由此可见，该事件不是孤立的年表点，而是资源、信息与权力关系重新配置的一环。不同地区的官员、社群和家庭可能以不同方式承担征发、保护、迁移或宗教活动；现存来源只能照亮其中一部分，不能替沉默者制造统一的经验。

\eventanchor{TK-LEDGER-004-1}{黄初三年（222年）·“曹丕册封孙权为吴王，试图以册命、质子要求和荆州牧名义约束长江下游政权”的前期动员与资源准备}账本记录曹魏与孙权集团在黄初三年（222年）发生““曹丕册封孙权为吴王，试图以册命、质子要求和荆州牧名义约束长江下游政权”的前期动员与资源准备”。条目把背景概括为当时的权力、资源与信息约束推动相关过程展开，并指出这一过程为“曹丕册封孙权为吴王，试图以册命、质子要求和荆州牧名义约束长江下游政权”提供可观察的条件。这个节点进入区域社会时，要经由文书、道路、军队、市场、寺院和家庭等不同中介；因此，政治行动的宣布、地方接收和日常生活的改变不能被视作同一时刻完成。

本条依据《三国志·吴主传》；《三国志·文帝纪》；《资治通鉴》魏纪，证据提示为“核心事项已有既有账本来源；本分解节点的参与者、执行范围和时间先后仍须逐案核验”。这些材料可支持年序、官方决定、战事或特定地点的线索，却难以直接统计所有人的财富、迁徙、信仰和动机。更稳妥的读法是区分诏令与实施、军报与人口、行政称谓与自我认同，并让地方材料的缺环限制解释的范围。

由此可见，该事件不是孤立的年表点，而是资源、信息与权力关系重新配置的一环。不同地区的官员、社群和家庭可能以不同方式承担征发、保护、迁移或宗教活动；现存来源只能照亮其中一部分，不能替沉默者制造统一的经验。

\eventanchor{TK-LEDGER-004-2}{黄初三年（222年）·“曹丕册封孙权为吴王，试图以册命、质子要求和荆州牧名义约束长江下游政权”的命令、联盟或地方响应}账本记录曹魏与孙权集团在黄初三年（222年）发生““曹丕册封孙权为吴王，试图以册命、质子要求和荆州牧名义约束长江下游政权”的命令、联盟或地方响应”。条目把背景概括为当时的权力、资源与信息约束推动相关过程展开，并指出这一过程为“曹丕册封孙权为吴王，试图以册命、质子要求和荆州牧名义约束长江下游政权”提供可观察的条件。这个节点进入区域社会时，要经由文书、道路、军队、市场、寺院和家庭等不同中介；因此，政治行动的宣布、地方接收和日常生活的改变不能被视作同一时刻完成。

本条依据《三国志·吴主传》；《三国志·文帝纪》；《资治通鉴》魏纪，证据提示为“核心事项已有既有账本来源；本分解节点的参与者、执行范围和时间先后仍须逐案核验”。这些材料可支持年序、官方决定、战事或特定地点的线索，却难以直接统计所有人的财富、迁徙、信仰和动机。更稳妥的读法是区分诏令与实施、军报与人口、行政称谓与自我认同，并让地方材料的缺环限制解释的范围。

由此可见，该事件不是孤立的年表点，而是资源、信息与权力关系重新配置的一环。不同地区的官员、社群和家庭可能以不同方式承担征发、保护、迁移或宗教活动；现存来源只能照亮其中一部分，不能替沉默者制造统一的经验。

\eventanchor{TK-LEDGER-004-3}{黄初三年（222年）·曹丕册封孙权为吴王，试图以册命、质子要求和荆州牧名义约束长江下游政权}账本记录曹魏与孙权集团在黄初三年（222年）发生“曹丕册封孙权为吴王，试图以册命、质子要求和荆州牧名义约束长江下游政权”。条目把背景概括为“曹丕册封孙权为吴王，试图以册命、质子要求和荆州牧名义约束长江下游政权”发生的政治与区域条件共同作用，并指出该节点改变后续资源、联盟或制度处置的条件。这个节点进入区域社会时，要经由文书、道路、军队、市场、寺院和家庭等不同中介；因此，政治行动的宣布、地方接收和日常生活的改变不能被视作同一时刻完成。

本条依据《三国志·吴主传》；《三国志·文帝纪》；《资治通鉴》魏纪，证据提示为“核心事项已有既有账本来源；本分解节点的参与者、执行范围和时间先后仍须逐案核验”。这些材料可支持年序、官方决定、战事或特定地点的线索，却难以直接统计所有人的财富、迁徙、信仰和动机。更稳妥的读法是区分诏令与实施、军报与人口、行政称谓与自我认同，并让地方材料的缺环限制解释的范围。

由此可见，该事件不是孤立的年表点，而是资源、信息与权力关系重新配置的一环。不同地区的官员、社群和家庭可能以不同方式承担征发、保护、迁移或宗教活动；现存来源只能照亮其中一部分，不能替沉默者制造统一的经验。

\eventanchor{TK-LEDGER-004-4}{黄初三年（222年）·“曹丕册封孙权为吴王，试图以册命、质子要求和荆州牧名义约束长江下游政权”的执行中的空间与人员处置}账本记录曹魏与孙权集团在黄初三年（222年）发生““曹丕册封孙权为吴王，试图以册命、质子要求和荆州牧名义约束长江下游政权”的执行中的空间与人员处置”。条目把背景概括为“曹丕册封孙权为吴王，试图以册命、质子要求和荆州牧名义约束长江下游政权”发生的政治与区域条件共同作用，并指出该节点改变后续资源、联盟或制度处置的条件。这个节点进入区域社会时，要经由文书、道路、军队、市场、寺院和家庭等不同中介；因此，政治行动的宣布、地方接收和日常生活的改变不能被视作同一时刻完成。

本条依据《三国志·吴主传》；《三国志·文帝纪》；《资治通鉴》魏纪，证据提示为“核心事项已有既有账本来源；本分解节点的参与者、执行范围和时间先后仍须逐案核验”。这些材料可支持年序、官方决定、战事或特定地点的线索，却难以直接统计所有人的财富、迁徙、信仰和动机。更稳妥的读法是区分诏令与实施、军报与人口、行政称谓与自我认同，并让地方材料的缺环限制解释的范围。

由此可见，该事件不是孤立的年表点，而是资源、信息与权力关系重新配置的一环。不同地区的官员、社群和家庭可能以不同方式承担征发、保护、迁移或宗教活动；现存来源只能照亮其中一部分，不能替沉默者制造统一的经验。

\eventanchor{TK-LEDGER-004-5}{黄初三年（222年）·“曹丕册封孙权为吴王，试图以册命、质子要求和荆州牧名义约束长江下游政权”的后续制度或军政调整}账本记录曹魏与孙权集团在黄初三年（222年）发生““曹丕册封孙权为吴王，试图以册命、质子要求和荆州牧名义约束长江下游政权”的后续制度或军政调整”。条目把背景概括为核心节点发生后仍须处理其遗留问题，并指出后续影响须在不同地区和材料类型中分别核验。这个节点进入区域社会时，要经由文书、道路、军队、市场、寺院和家庭等不同中介；因此，政治行动的宣布、地方接收和日常生活的改变不能被视作同一时刻完成。

本条依据《三国志·吴主传》；《三国志·文帝纪》；《资治通鉴》魏纪，证据提示为“核心事项已有既有账本来源；本分解节点的参与者、执行范围和时间先后仍须逐案核验”。这些材料可支持年序、官方决定、战事或特定地点的线索，却难以直接统计所有人的财富、迁徙、信仰和动机。更稳妥的读法是区分诏令与实施、军报与人口、行政称谓与自我认同，并让地方材料的缺环限制解释的范围。

由此可见，该事件不是孤立的年表点，而是资源、信息与权力关系重新配置的一环。不同地区的官员、社群和家庭可能以不同方式承担征发、保护、迁移或宗教活动；现存来源只能照亮其中一部分，不能替沉默者制造统一的经验。

\eventanchor{TK-LEDGER-004-6}{黄初三年（222年）·“曹丕册封孙权为吴王，试图以册命、质子要求和荆州牧名义约束长江下游政权”的异源材料与影响范围的复核}账本记录曹魏与孙权集团在黄初三年（222年）发生““曹丕册封孙权为吴王，试图以册命、质子要求和荆州牧名义约束长江下游政权”的异源材料与影响范围的复核”。条目把背景概括为核心节点发生后仍须处理其遗留问题，并指出后续影响须在不同地区和材料类型中分别核验。这个节点进入区域社会时，要经由文书、道路、军队、市场、寺院和家庭等不同中介；因此，政治行动的宣布、地方接收和日常生活的改变不能被视作同一时刻完成。

本条依据《三国志·吴主传》；《三国志·文帝纪》；《资治通鉴》魏纪，证据提示为“核心事项已有既有账本来源；本分解节点的参与者、执行范围和时间先后仍须逐案核验”。这些材料可支持年序、官方决定、战事或特定地点的线索，却难以直接统计所有人的财富、迁徙、信仰和动机。更稳妥的读法是区分诏令与实施、军报与人口、行政称谓与自我认同，并让地方材料的缺环限制解释的范围。

由此可见，该事件不是孤立的年表点，而是资源、信息与权力关系重新配置的一环。不同地区的官员、社群和家庭可能以不同方式承担征发、保护、迁移或宗教活动；现存来源只能照亮其中一部分，不能替沉默者制造统一的经验。

\eventanchor{TK-LEDGER-005-1}{章武二年（222年）·“陆逊统吴军在夷陵、猇亭一线击破刘备，蜀军沿长江退向白帝”的前期动员与资源准备}账本记录蜀汉与孙吴在章武二年（222年）发生““陆逊统吴军在夷陵、猇亭一线击破刘备，蜀军沿长江退向白帝”的前期动员与资源准备”。条目把背景概括为当时的权力、资源与信息约束推动相关过程展开，并指出这一过程为“陆逊统吴军在夷陵、猇亭一线击破刘备，蜀军沿长江退向白帝”提供可观察的条件。这个节点进入区域社会时，要经由文书、道路、军队、市场、寺院和家庭等不同中介；因此，政治行动的宣布、地方接收和日常生活的改变不能被视作同一时刻完成。

本条依据《三国志·先主传》；《三国志·陆逊传》；《三国志·吴主传》；《资治通鉴》魏纪，证据提示为“核心事项已有既有账本来源；本分解节点的参与者、执行范围和时间先后仍须逐案核验”。这些材料可支持年序、官方决定、战事或特定地点的线索，却难以直接统计所有人的财富、迁徙、信仰和动机。更稳妥的读法是区分诏令与实施、军报与人口、行政称谓与自我认同，并让地方材料的缺环限制解释的范围。

由此可见，该事件不是孤立的年表点，而是资源、信息与权力关系重新配置的一环。不同地区的官员、社群和家庭可能以不同方式承担征发、保护、迁移或宗教活动；现存来源只能照亮其中一部分，不能替沉默者制造统一的经验。

\eventanchor{TK-LEDGER-005-2}{章武二年（222年）·“陆逊统吴军在夷陵、猇亭一线击破刘备，蜀军沿长江退向白帝”的命令、联盟或地方响应}账本记录蜀汉与孙吴在章武二年（222年）发生““陆逊统吴军在夷陵、猇亭一线击破刘备，蜀军沿长江退向白帝”的命令、联盟或地方响应”。条目把背景概括为当时的权力、资源与信息约束推动相关过程展开，并指出这一过程为“陆逊统吴军在夷陵、猇亭一线击破刘备，蜀军沿长江退向白帝”提供可观察的条件。这个节点进入区域社会时，要经由文书、道路、军队、市场、寺院和家庭等不同中介；因此，政治行动的宣布、地方接收和日常生活的改变不能被视作同一时刻完成。

本条依据《三国志·先主传》；《三国志·陆逊传》；《三国志·吴主传》；《资治通鉴》魏纪，证据提示为“核心事项已有既有账本来源；本分解节点的参与者、执行范围和时间先后仍须逐案核验”。这些材料可支持年序、官方决定、战事或特定地点的线索，却难以直接统计所有人的财富、迁徙、信仰和动机。更稳妥的读法是区分诏令与实施、军报与人口、行政称谓与自我认同，并让地方材料的缺环限制解释的范围。

由此可见，该事件不是孤立的年表点，而是资源、信息与权力关系重新配置的一环。不同地区的官员、社群和家庭可能以不同方式承担征发、保护、迁移或宗教活动；现存来源只能照亮其中一部分，不能替沉默者制造统一的经验。

\eventanchor{TK-LEDGER-005-3}{章武二年（222年）·陆逊统吴军在夷陵、猇亭一线击破刘备，蜀军沿长江退向白帝}账本记录蜀汉与孙吴在章武二年（222年）发生“陆逊统吴军在夷陵、猇亭一线击破刘备，蜀军沿长江退向白帝”。条目把背景概括为“陆逊统吴军在夷陵、猇亭一线击破刘备，蜀军沿长江退向白帝”发生的政治与区域条件共同作用，并指出该节点改变后续资源、联盟或制度处置的条件。这个节点进入区域社会时，要经由文书、道路、军队、市场、寺院和家庭等不同中介；因此，政治行动的宣布、地方接收和日常生活的改变不能被视作同一时刻完成。

本条依据《三国志·先主传》；《三国志·陆逊传》；《三国志·吴主传》；《资治通鉴》魏纪，证据提示为“核心事项已有既有账本来源；本分解节点的参与者、执行范围和时间先后仍须逐案核验”。这些材料可支持年序、官方决定、战事或特定地点的线索，却难以直接统计所有人的财富、迁徙、信仰和动机。更稳妥的读法是区分诏令与实施、军报与人口、行政称谓与自我认同，并让地方材料的缺环限制解释的范围。

由此可见，该事件不是孤立的年表点，而是资源、信息与权力关系重新配置的一环。不同地区的官员、社群和家庭可能以不同方式承担征发、保护、迁移或宗教活动；现存来源只能照亮其中一部分，不能替沉默者制造统一的经验。

\eventanchor{TK-LEDGER-005-4}{章武二年（222年）·“陆逊统吴军在夷陵、猇亭一线击破刘备，蜀军沿长江退向白帝”的执行中的空间与人员处置}账本记录蜀汉与孙吴在章武二年（222年）发生““陆逊统吴军在夷陵、猇亭一线击破刘备，蜀军沿长江退向白帝”的执行中的空间与人员处置”。条目把背景概括为“陆逊统吴军在夷陵、猇亭一线击破刘备，蜀军沿长江退向白帝”发生的政治与区域条件共同作用，并指出该节点改变后续资源、联盟或制度处置的条件。这个节点进入区域社会时，要经由文书、道路、军队、市场、寺院和家庭等不同中介；因此，政治行动的宣布、地方接收和日常生活的改变不能被视作同一时刻完成。

本条依据《三国志·先主传》；《三国志·陆逊传》；《三国志·吴主传》；《资治通鉴》魏纪，证据提示为“核心事项已有既有账本来源；本分解节点的参与者、执行范围和时间先后仍须逐案核验”。这些材料可支持年序、官方决定、战事或特定地点的线索，却难以直接统计所有人的财富、迁徙、信仰和动机。更稳妥的读法是区分诏令与实施、军报与人口、行政称谓与自我认同，并让地方材料的缺环限制解释的范围。

由此可见，该事件不是孤立的年表点，而是资源、信息与权力关系重新配置的一环。不同地区的官员、社群和家庭可能以不同方式承担征发、保护、迁移或宗教活动；现存来源只能照亮其中一部分，不能替沉默者制造统一的经验。

\eventanchor{TK-LEDGER-005-5}{章武二年（222年）·“陆逊统吴军在夷陵、猇亭一线击破刘备，蜀军沿长江退向白帝”的后续制度或军政调整}账本记录蜀汉与孙吴在章武二年（222年）发生““陆逊统吴军在夷陵、猇亭一线击破刘备，蜀军沿长江退向白帝”的后续制度或军政调整”。条目把背景概括为核心节点发生后仍须处理其遗留问题，并指出后续影响须在不同地区和材料类型中分别核验。这个节点进入区域社会时，要经由文书、道路、军队、市场、寺院和家庭等不同中介；因此，政治行动的宣布、地方接收和日常生活的改变不能被视作同一时刻完成。

本条依据《三国志·先主传》；《三国志·陆逊传》；《三国志·吴主传》；《资治通鉴》魏纪，证据提示为“核心事项已有既有账本来源；本分解节点的参与者、执行范围和时间先后仍须逐案核验”。这些材料可支持年序、官方决定、战事或特定地点的线索，却难以直接统计所有人的财富、迁徙、信仰和动机。更稳妥的读法是区分诏令与实施、军报与人口、行政称谓与自我认同，并让地方材料的缺环限制解释的范围。

由此可见，该事件不是孤立的年表点，而是资源、信息与权力关系重新配置的一环。不同地区的官员、社群和家庭可能以不同方式承担征发、保护、迁移或宗教活动；现存来源只能照亮其中一部分，不能替沉默者制造统一的经验。

\eventanchor{TK-LEDGER-005-6}{章武二年（222年）·“陆逊统吴军在夷陵、猇亭一线击破刘备，蜀军沿长江退向白帝”的异源材料与影响范围的复核}账本记录蜀汉与孙吴在章武二年（222年）发生““陆逊统吴军在夷陵、猇亭一线击破刘备，蜀军沿长江退向白帝”的异源材料与影响范围的复核”。条目把背景概括为核心节点发生后仍须处理其遗留问题，并指出后续影响须在不同地区和材料类型中分别核验。这个节点进入区域社会时，要经由文书、道路、军队、市场、寺院和家庭等不同中介；因此，政治行动的宣布、地方接收和日常生活的改变不能被视作同一时刻完成。

本条依据《三国志·先主传》；《三国志·陆逊传》；《三国志·吴主传》；《资治通鉴》魏纪，证据提示为“核心事项已有既有账本来源；本分解节点的参与者、执行范围和时间先后仍须逐案核验”。这些材料可支持年序、官方决定、战事或特定地点的线索，却难以直接统计所有人的财富、迁徙、信仰和动机。更稳妥的读法是区分诏令与实施、军报与人口、行政称谓与自我认同，并让地方材料的缺环限制解释的范围。

由此可见，该事件不是孤立的年表点，而是资源、信息与权力关系重新配置的一环。不同地区的官员、社群和家庭可能以不同方式承担征发、保护、迁移或宗教活动；现存来源只能照亮其中一部分，不能替沉默者制造统一的经验。

\subsection{社会关系与宗教空间}

\eventanchor{TK-LEDGER-006-1}{章武三年四月（223年）·“刘备在永安去世，刘禅继位，诸葛亮受遗诏辅政”的前期动员与资源准备}账本记录蜀汉在章武三年四月（223年）发生““刘备在永安去世，刘禅继位，诸葛亮受遗诏辅政”的前期动员与资源准备”。条目把背景概括为当时的权力、资源与信息约束推动相关过程展开，并指出这一过程为“刘备在永安去世，刘禅继位，诸葛亮受遗诏辅政”提供可观察的条件。这个节点进入区域社会时，要经由文书、道路、军队、市场、寺院和家庭等不同中介；因此，政治行动的宣布、地方接收和日常生活的改变不能被视作同一时刻完成。

本条依据《三国志·先主传》；《三国志·后主传》；《三国志·诸葛亮传》，证据提示为“核心事项已有既有账本来源；本分解节点的参与者、执行范围和时间先后仍须逐案核验”。这些材料可支持年序、官方决定、战事或特定地点的线索，却难以直接统计所有人的财富、迁徙、信仰和动机。更稳妥的读法是区分诏令与实施、军报与人口、行政称谓与自我认同，并让地方材料的缺环限制解释的范围。

由此可见，该事件不是孤立的年表点，而是资源、信息与权力关系重新配置的一环。不同地区的官员、社群和家庭可能以不同方式承担征发、保护、迁移或宗教活动；现存来源只能照亮其中一部分，不能替沉默者制造统一的经验。

\eventanchor{TK-LEDGER-006-2}{章武三年四月（223年）·“刘备在永安去世，刘禅继位，诸葛亮受遗诏辅政”的命令、联盟或地方响应}账本记录蜀汉在章武三年四月（223年）发生““刘备在永安去世，刘禅继位，诸葛亮受遗诏辅政”的命令、联盟或地方响应”。条目把背景概括为当时的权力、资源与信息约束推动相关过程展开，并指出这一过程为“刘备在永安去世，刘禅继位，诸葛亮受遗诏辅政”提供可观察的条件。这个节点进入区域社会时，要经由文书、道路、军队、市场、寺院和家庭等不同中介；因此，政治行动的宣布、地方接收和日常生活的改变不能被视作同一时刻完成。

本条依据《三国志·先主传》；《三国志·后主传》；《三国志·诸葛亮传》，证据提示为“核心事项已有既有账本来源；本分解节点的参与者、执行范围和时间先后仍须逐案核验”。这些材料可支持年序、官方决定、战事或特定地点的线索，却难以直接统计所有人的财富、迁徙、信仰和动机。更稳妥的读法是区分诏令与实施、军报与人口、行政称谓与自我认同，并让地方材料的缺环限制解释的范围。

由此可见，该事件不是孤立的年表点，而是资源、信息与权力关系重新配置的一环。不同地区的官员、社群和家庭可能以不同方式承担征发、保护、迁移或宗教活动；现存来源只能照亮其中一部分，不能替沉默者制造统一的经验。

\eventanchor{TK-LEDGER-006-3}{章武三年四月（223年）·刘备在永安去世，刘禅继位，诸葛亮受遗诏辅政}账本记录蜀汉在章武三年四月（223年）发生“刘备在永安去世，刘禅继位，诸葛亮受遗诏辅政”。条目把背景概括为“刘备在永安去世，刘禅继位，诸葛亮受遗诏辅政”发生的政治与区域条件共同作用，并指出该节点改变后续资源、联盟或制度处置的条件。这个节点进入区域社会时，要经由文书、道路、军队、市场、寺院和家庭等不同中介；因此，政治行动的宣布、地方接收和日常生活的改变不能被视作同一时刻完成。

本条依据《三国志·先主传》；《三国志·后主传》；《三国志·诸葛亮传》，证据提示为“核心事项已有既有账本来源；本分解节点的参与者、执行范围和时间先后仍须逐案核验”。这些材料可支持年序、官方决定、战事或特定地点的线索，却难以直接统计所有人的财富、迁徙、信仰和动机。更稳妥的读法是区分诏令与实施、军报与人口、行政称谓与自我认同，并让地方材料的缺环限制解释的范围。

由此可见，该事件不是孤立的年表点，而是资源、信息与权力关系重新配置的一环。不同地区的官员、社群和家庭可能以不同方式承担征发、保护、迁移或宗教活动；现存来源只能照亮其中一部分，不能替沉默者制造统一的经验。

\eventanchor{TK-LEDGER-006-4}{章武三年四月（223年）·“刘备在永安去世，刘禅继位，诸葛亮受遗诏辅政”的执行中的空间与人员处置}账本记录蜀汉在章武三年四月（223年）发生““刘备在永安去世，刘禅继位，诸葛亮受遗诏辅政”的执行中的空间与人员处置”。条目把背景概括为“刘备在永安去世，刘禅继位，诸葛亮受遗诏辅政”发生的政治与区域条件共同作用，并指出该节点改变后续资源、联盟或制度处置的条件。这个节点进入区域社会时，要经由文书、道路、军队、市场、寺院和家庭等不同中介；因此，政治行动的宣布、地方接收和日常生活的改变不能被视作同一时刻完成。

本条依据《三国志·先主传》；《三国志·后主传》；《三国志·诸葛亮传》，证据提示为“核心事项已有既有账本来源；本分解节点的参与者、执行范围和时间先后仍须逐案核验”。这些材料可支持年序、官方决定、战事或特定地点的线索，却难以直接统计所有人的财富、迁徙、信仰和动机。更稳妥的读法是区分诏令与实施、军报与人口、行政称谓与自我认同，并让地方材料的缺环限制解释的范围。

由此可见，该事件不是孤立的年表点，而是资源、信息与权力关系重新配置的一环。不同地区的官员、社群和家庭可能以不同方式承担征发、保护、迁移或宗教活动；现存来源只能照亮其中一部分，不能替沉默者制造统一的经验。

\eventanchor{TK-LEDGER-006-5}{章武三年四月（223年）·“刘备在永安去世，刘禅继位，诸葛亮受遗诏辅政”的后续制度或军政调整}账本记录蜀汉在章武三年四月（223年）发生““刘备在永安去世，刘禅继位，诸葛亮受遗诏辅政”的后续制度或军政调整”。条目把背景概括为核心节点发生后仍须处理其遗留问题，并指出后续影响须在不同地区和材料类型中分别核验。这个节点进入区域社会时，要经由文书、道路、军队、市场、寺院和家庭等不同中介；因此，政治行动的宣布、地方接收和日常生活的改变不能被视作同一时刻完成。

本条依据《三国志·先主传》；《三国志·后主传》；《三国志·诸葛亮传》，证据提示为“核心事项已有既有账本来源；本分解节点的参与者、执行范围和时间先后仍须逐案核验”。这些材料可支持年序、官方决定、战事或特定地点的线索，却难以直接统计所有人的财富、迁徙、信仰和动机。更稳妥的读法是区分诏令与实施、军报与人口、行政称谓与自我认同，并让地方材料的缺环限制解释的范围。

由此可见，该事件不是孤立的年表点，而是资源、信息与权力关系重新配置的一环。不同地区的官员、社群和家庭可能以不同方式承担征发、保护、迁移或宗教活动；现存来源只能照亮其中一部分，不能替沉默者制造统一的经验。

\eventanchor{TK-LEDGER-006-6}{章武三年四月（223年）·“刘备在永安去世，刘禅继位，诸葛亮受遗诏辅政”的异源材料与影响范围的复核}账本记录蜀汉在章武三年四月（223年）发生““刘备在永安去世，刘禅继位，诸葛亮受遗诏辅政”的异源材料与影响范围的复核”。条目把背景概括为核心节点发生后仍须处理其遗留问题，并指出后续影响须在不同地区和材料类型中分别核验。这个节点进入区域社会时，要经由文书、道路、军队、市场、寺院和家庭等不同中介；因此，政治行动的宣布、地方接收和日常生活的改变不能被视作同一时刻完成。

本条依据《三国志·先主传》；《三国志·后主传》；《三国志·诸葛亮传》，证据提示为“核心事项已有既有账本来源；本分解节点的参与者、执行范围和时间先后仍须逐案核验”。这些材料可支持年序、官方决定、战事或特定地点的线索，却难以直接统计所有人的财富、迁徙、信仰和动机。更稳妥的读法是区分诏令与实施、军报与人口、行政称谓与自我认同，并让地方材料的缺环限制解释的范围。

由此可见，该事件不是孤立的年表点，而是资源、信息与权力关系重新配置的一环。不同地区的官员、社群和家庭可能以不同方式承担征发、保护、迁移或宗教活动；现存来源只能照亮其中一部分，不能替沉默者制造统一的经验。

\eventanchor{TK-LEDGER-007-1}{建兴元年（223年）·“邓芝奉使孙吴，孙权重新与蜀汉通好，双方恢复共同抗魏的外交框架”的前期动员与资源准备}账本记录蜀汉与孙吴在建兴元年（223年）发生““邓芝奉使孙吴，孙权重新与蜀汉通好，双方恢复共同抗魏的外交框架”的前期动员与资源准备”。条目把背景概括为当时的权力、资源与信息约束推动相关过程展开，并指出这一过程为“邓芝奉使孙吴，孙权重新与蜀汉通好，双方恢复共同抗魏的外交框架”提供可观察的条件。这个节点进入区域社会时，要经由文书、道路、军队、市场、寺院和家庭等不同中介；因此，政治行动的宣布、地方接收和日常生活的改变不能被视作同一时刻完成。

本条依据《三国志·邓芝传》；《三国志·吴主传》；《三国志·诸葛亮传》，证据提示为“核心事项已有既有账本来源；本分解节点的参与者、执行范围和时间先后仍须逐案核验”。这些材料可支持年序、官方决定、战事或特定地点的线索，却难以直接统计所有人的财富、迁徙、信仰和动机。更稳妥的读法是区分诏令与实施、军报与人口、行政称谓与自我认同，并让地方材料的缺环限制解释的范围。

由此可见，该事件不是孤立的年表点，而是资源、信息与权力关系重新配置的一环。不同地区的官员、社群和家庭可能以不同方式承担征发、保护、迁移或宗教活动；现存来源只能照亮其中一部分，不能替沉默者制造统一的经验。

\eventanchor{TK-LEDGER-007-2}{建兴元年（223年）·“邓芝奉使孙吴，孙权重新与蜀汉通好，双方恢复共同抗魏的外交框架”的命令、联盟或地方响应}账本记录蜀汉与孙吴在建兴元年（223年）发生““邓芝奉使孙吴，孙权重新与蜀汉通好，双方恢复共同抗魏的外交框架”的命令、联盟或地方响应”。条目把背景概括为当时的权力、资源与信息约束推动相关过程展开，并指出这一过程为“邓芝奉使孙吴，孙权重新与蜀汉通好，双方恢复共同抗魏的外交框架”提供可观察的条件。这个节点进入区域社会时，要经由文书、道路、军队、市场、寺院和家庭等不同中介；因此，政治行动的宣布、地方接收和日常生活的改变不能被视作同一时刻完成。

本条依据《三国志·邓芝传》；《三国志·吴主传》；《三国志·诸葛亮传》，证据提示为“核心事项已有既有账本来源；本分解节点的参与者、执行范围和时间先后仍须逐案核验”。这些材料可支持年序、官方决定、战事或特定地点的线索，却难以直接统计所有人的财富、迁徙、信仰和动机。更稳妥的读法是区分诏令与实施、军报与人口、行政称谓与自我认同，并让地方材料的缺环限制解释的范围。

由此可见，该事件不是孤立的年表点，而是资源、信息与权力关系重新配置的一环。不同地区的官员、社群和家庭可能以不同方式承担征发、保护、迁移或宗教活动；现存来源只能照亮其中一部分，不能替沉默者制造统一的经验。

\eventanchor{TK-LEDGER-007-3}{建兴元年（223年）·邓芝奉使孙吴，孙权重新与蜀汉通好，双方恢复共同抗魏的外交框架}账本记录蜀汉与孙吴在建兴元年（223年）发生“邓芝奉使孙吴，孙权重新与蜀汉通好，双方恢复共同抗魏的外交框架”。条目把背景概括为“邓芝奉使孙吴，孙权重新与蜀汉通好，双方恢复共同抗魏的外交框架”发生的政治与区域条件共同作用，并指出该节点改变后续资源、联盟或制度处置的条件。这个节点进入区域社会时，要经由文书、道路、军队、市场、寺院和家庭等不同中介；因此，政治行动的宣布、地方接收和日常生活的改变不能被视作同一时刻完成。

本条依据《三国志·邓芝传》；《三国志·吴主传》；《三国志·诸葛亮传》，证据提示为“核心事项已有既有账本来源；本分解节点的参与者、执行范围和时间先后仍须逐案核验”。这些材料可支持年序、官方决定、战事或特定地点的线索，却难以直接统计所有人的财富、迁徙、信仰和动机。更稳妥的读法是区分诏令与实施、军报与人口、行政称谓与自我认同，并让地方材料的缺环限制解释的范围。

由此可见，该事件不是孤立的年表点，而是资源、信息与权力关系重新配置的一环。不同地区的官员、社群和家庭可能以不同方式承担征发、保护、迁移或宗教活动；现存来源只能照亮其中一部分，不能替沉默者制造统一的经验。

\eventanchor{TK-LEDGER-007-4}{建兴元年（223年）·“邓芝奉使孙吴，孙权重新与蜀汉通好，双方恢复共同抗魏的外交框架”的执行中的空间与人员处置}账本记录蜀汉与孙吴在建兴元年（223年）发生““邓芝奉使孙吴，孙权重新与蜀汉通好，双方恢复共同抗魏的外交框架”的执行中的空间与人员处置”。条目把背景概括为“邓芝奉使孙吴，孙权重新与蜀汉通好，双方恢复共同抗魏的外交框架”发生的政治与区域条件共同作用，并指出该节点改变后续资源、联盟或制度处置的条件。这个节点进入区域社会时，要经由文书、道路、军队、市场、寺院和家庭等不同中介；因此，政治行动的宣布、地方接收和日常生活的改变不能被视作同一时刻完成。

本条依据《三国志·邓芝传》；《三国志·吴主传》；《三国志·诸葛亮传》，证据提示为“核心事项已有既有账本来源；本分解节点的参与者、执行范围和时间先后仍须逐案核验”。这些材料可支持年序、官方决定、战事或特定地点的线索，却难以直接统计所有人的财富、迁徙、信仰和动机。更稳妥的读法是区分诏令与实施、军报与人口、行政称谓与自我认同，并让地方材料的缺环限制解释的范围。

由此可见，该事件不是孤立的年表点，而是资源、信息与权力关系重新配置的一环。不同地区的官员、社群和家庭可能以不同方式承担征发、保护、迁移或宗教活动；现存来源只能照亮其中一部分，不能替沉默者制造统一的经验。

\eventanchor{TK-LEDGER-007-5}{建兴元年（223年）·“邓芝奉使孙吴，孙权重新与蜀汉通好，双方恢复共同抗魏的外交框架”的后续制度或军政调整}账本记录蜀汉与孙吴在建兴元年（223年）发生““邓芝奉使孙吴，孙权重新与蜀汉通好，双方恢复共同抗魏的外交框架”的后续制度或军政调整”。条目把背景概括为核心节点发生后仍须处理其遗留问题，并指出后续影响须在不同地区和材料类型中分别核验。这个节点进入区域社会时，要经由文书、道路、军队、市场、寺院和家庭等不同中介；因此，政治行动的宣布、地方接收和日常生活的改变不能被视作同一时刻完成。

本条依据《三国志·邓芝传》；《三国志·吴主传》；《三国志·诸葛亮传》，证据提示为“核心事项已有既有账本来源；本分解节点的参与者、执行范围和时间先后仍须逐案核验”。这些材料可支持年序、官方决定、战事或特定地点的线索，却难以直接统计所有人的财富、迁徙、信仰和动机。更稳妥的读法是区分诏令与实施、军报与人口、行政称谓与自我认同，并让地方材料的缺环限制解释的范围。

由此可见，该事件不是孤立的年表点，而是资源、信息与权力关系重新配置的一环。不同地区的官员、社群和家庭可能以不同方式承担征发、保护、迁移或宗教活动；现存来源只能照亮其中一部分，不能替沉默者制造统一的经验。

\eventanchor{TK-LEDGER-007-6}{建兴元年（223年）·“邓芝奉使孙吴，孙权重新与蜀汉通好，双方恢复共同抗魏的外交框架”的异源材料与影响范围的复核}账本记录蜀汉与孙吴在建兴元年（223年）发生““邓芝奉使孙吴，孙权重新与蜀汉通好，双方恢复共同抗魏的外交框架”的异源材料与影响范围的复核”。条目把背景概括为核心节点发生后仍须处理其遗留问题，并指出后续影响须在不同地区和材料类型中分别核验。这个节点进入区域社会时，要经由文书、道路、军队、市场、寺院和家庭等不同中介；因此，政治行动的宣布、地方接收和日常生活的改变不能被视作同一时刻完成。

本条依据《三国志·邓芝传》；《三国志·吴主传》；《三国志·诸葛亮传》，证据提示为“核心事项已有既有账本来源；本分解节点的参与者、执行范围和时间先后仍须逐案核验”。这些材料可支持年序、官方决定、战事或特定地点的线索，却难以直接统计所有人的财富、迁徙、信仰和动机。更稳妥的读法是区分诏令与实施、军报与人口、行政称谓与自我认同，并让地方材料的缺环限制解释的范围。

由此可见，该事件不是孤立的年表点，而是资源、信息与权力关系重新配置的一环。不同地区的官员、社群和家庭可能以不同方式承担征发、保护、迁移或宗教活动；现存来源只能照亮其中一部分，不能替沉默者制造统一的经验。

\eventanchor{TK-LEDGER-008-1}{黄初五年（224年）·“曹丕至广陵观察南征，因长江冰冻、舟师与渡江条件不具而撤军”的前期动员与资源准备}账本记录曹魏与孙吴在黄初五年（224年）发生““曹丕至广陵观察南征，因长江冰冻、舟师与渡江条件不具而撤军”的前期动员与资源准备”。条目把背景概括为当时的权力、资源与信息约束推动相关过程展开，并指出这一过程为“曹丕至广陵观察南征，因长江冰冻、舟师与渡江条件不具而撤军”提供可观察的条件。这个节点进入区域社会时，要经由文书、道路、军队、市场、寺院和家庭等不同中介；因此，政治行动的宣布、地方接收和日常生活的改变不能被视作同一时刻完成。

本条依据《三国志·文帝纪》；《三国志·吴主传》；《资治通鉴》魏纪，证据提示为“核心事项已有既有账本来源；本分解节点的参与者、执行范围和时间先后仍须逐案核验”。这些材料可支持年序、官方决定、战事或特定地点的线索，却难以直接统计所有人的财富、迁徙、信仰和动机。更稳妥的读法是区分诏令与实施、军报与人口、行政称谓与自我认同，并让地方材料的缺环限制解释的范围。

由此可见，该事件不是孤立的年表点，而是资源、信息与权力关系重新配置的一环。不同地区的官员、社群和家庭可能以不同方式承担征发、保护、迁移或宗教活动；现存来源只能照亮其中一部分，不能替沉默者制造统一的经验。

\eventanchor{TK-LEDGER-008-2}{黄初五年（224年）·“曹丕至广陵观察南征，因长江冰冻、舟师与渡江条件不具而撤军”的命令、联盟或地方响应}账本记录曹魏与孙吴在黄初五年（224年）发生““曹丕至广陵观察南征，因长江冰冻、舟师与渡江条件不具而撤军”的命令、联盟或地方响应”。条目把背景概括为当时的权力、资源与信息约束推动相关过程展开，并指出这一过程为“曹丕至广陵观察南征，因长江冰冻、舟师与渡江条件不具而撤军”提供可观察的条件。这个节点进入区域社会时，要经由文书、道路、军队、市场、寺院和家庭等不同中介；因此，政治行动的宣布、地方接收和日常生活的改变不能被视作同一时刻完成。

本条依据《三国志·文帝纪》；《三国志·吴主传》；《资治通鉴》魏纪，证据提示为“核心事项已有既有账本来源；本分解节点的参与者、执行范围和时间先后仍须逐案核验”。这些材料可支持年序、官方决定、战事或特定地点的线索，却难以直接统计所有人的财富、迁徙、信仰和动机。更稳妥的读法是区分诏令与实施、军报与人口、行政称谓与自我认同，并让地方材料的缺环限制解释的范围。

由此可见，该事件不是孤立的年表点，而是资源、信息与权力关系重新配置的一环。不同地区的官员、社群和家庭可能以不同方式承担征发、保护、迁移或宗教活动；现存来源只能照亮其中一部分，不能替沉默者制造统一的经验。

\eventanchor{TK-LEDGER-008-3}{黄初五年（224年）·曹丕至广陵观察南征，因长江冰冻、舟师与渡江条件不具而撤军}账本记录曹魏与孙吴在黄初五年（224年）发生“曹丕至广陵观察南征，因长江冰冻、舟师与渡江条件不具而撤军”。条目把背景概括为“曹丕至广陵观察南征，因长江冰冻、舟师与渡江条件不具而撤军”发生的政治与区域条件共同作用，并指出该节点改变后续资源、联盟或制度处置的条件。这个节点进入区域社会时，要经由文书、道路、军队、市场、寺院和家庭等不同中介；因此，政治行动的宣布、地方接收和日常生活的改变不能被视作同一时刻完成。

本条依据《三国志·文帝纪》；《三国志·吴主传》；《资治通鉴》魏纪，证据提示为“核心事项已有既有账本来源；本分解节点的参与者、执行范围和时间先后仍须逐案核验”。这些材料可支持年序、官方决定、战事或特定地点的线索，却难以直接统计所有人的财富、迁徙、信仰和动机。更稳妥的读法是区分诏令与实施、军报与人口、行政称谓与自我认同，并让地方材料的缺环限制解释的范围。

由此可见，该事件不是孤立的年表点，而是资源、信息与权力关系重新配置的一环。不同地区的官员、社群和家庭可能以不同方式承担征发、保护、迁移或宗教活动；现存来源只能照亮其中一部分，不能替沉默者制造统一的经验。

\eventanchor{TK-LEDGER-008-4}{黄初五年（224年）·“曹丕至广陵观察南征，因长江冰冻、舟师与渡江条件不具而撤军”的执行中的空间与人员处置}账本记录曹魏与孙吴在黄初五年（224年）发生““曹丕至广陵观察南征，因长江冰冻、舟师与渡江条件不具而撤军”的执行中的空间与人员处置”。条目把背景概括为“曹丕至广陵观察南征，因长江冰冻、舟师与渡江条件不具而撤军”发生的政治与区域条件共同作用，并指出该节点改变后续资源、联盟或制度处置的条件。这个节点进入区域社会时，要经由文书、道路、军队、市场、寺院和家庭等不同中介；因此，政治行动的宣布、地方接收和日常生活的改变不能被视作同一时刻完成。

本条依据《三国志·文帝纪》；《三国志·吴主传》；《资治通鉴》魏纪，证据提示为“核心事项已有既有账本来源；本分解节点的参与者、执行范围和时间先后仍须逐案核验”。这些材料可支持年序、官方决定、战事或特定地点的线索，却难以直接统计所有人的财富、迁徙、信仰和动机。更稳妥的读法是区分诏令与实施、军报与人口、行政称谓与自我认同，并让地方材料的缺环限制解释的范围。

由此可见，该事件不是孤立的年表点，而是资源、信息与权力关系重新配置的一环。不同地区的官员、社群和家庭可能以不同方式承担征发、保护、迁移或宗教活动；现存来源只能照亮其中一部分，不能替沉默者制造统一的经验。

\eventanchor{TK-LEDGER-008-5}{黄初五年（224年）·“曹丕至广陵观察南征，因长江冰冻、舟师与渡江条件不具而撤军”的后续制度或军政调整}账本记录曹魏与孙吴在黄初五年（224年）发生““曹丕至广陵观察南征，因长江冰冻、舟师与渡江条件不具而撤军”的后续制度或军政调整”。条目把背景概括为核心节点发生后仍须处理其遗留问题，并指出后续影响须在不同地区和材料类型中分别核验。这个节点进入区域社会时，要经由文书、道路、军队、市场、寺院和家庭等不同中介；因此，政治行动的宣布、地方接收和日常生活的改变不能被视作同一时刻完成。

本条依据《三国志·文帝纪》；《三国志·吴主传》；《资治通鉴》魏纪，证据提示为“核心事项已有既有账本来源；本分解节点的参与者、执行范围和时间先后仍须逐案核验”。这些材料可支持年序、官方决定、战事或特定地点的线索，却难以直接统计所有人的财富、迁徙、信仰和动机。更稳妥的读法是区分诏令与实施、军报与人口、行政称谓与自我认同，并让地方材料的缺环限制解释的范围。

由此可见，该事件不是孤立的年表点，而是资源、信息与权力关系重新配置的一环。不同地区的官员、社群和家庭可能以不同方式承担征发、保护、迁移或宗教活动；现存来源只能照亮其中一部分，不能替沉默者制造统一的经验。

\eventanchor{TK-LEDGER-008-6}{黄初五年（224年）·“曹丕至广陵观察南征，因长江冰冻、舟师与渡江条件不具而撤军”的异源材料与影响范围的复核}账本记录曹魏与孙吴在黄初五年（224年）发生““曹丕至广陵观察南征，因长江冰冻、舟师与渡江条件不具而撤军”的异源材料与影响范围的复核”。条目把背景概括为核心节点发生后仍须处理其遗留问题，并指出后续影响须在不同地区和材料类型中分别核验。这个节点进入区域社会时，要经由文书、道路、军队、市场、寺院和家庭等不同中介；因此，政治行动的宣布、地方接收和日常生活的改变不能被视作同一时刻完成。

本条依据《三国志·文帝纪》；《三国志·吴主传》；《资治通鉴》魏纪，证据提示为“核心事项已有既有账本来源；本分解节点的参与者、执行范围和时间先后仍须逐案核验”。这些材料可支持年序、官方决定、战事或特定地点的线索，却难以直接统计所有人的财富、迁徙、信仰和动机。更稳妥的读法是区分诏令与实施、军报与人口、行政称谓与自我认同，并让地方材料的缺环限制解释的范围。

由此可见，该事件不是孤立的年表点，而是资源、信息与权力关系重新配置的一环。不同地区的官员、社群和家庭可能以不同方式承担征发、保护、迁移或宗教活动；现存来源只能照亮其中一部分，不能替沉默者制造统一的经验。

\eventanchor{TK-LEDGER-009-1}{建兴三年（225年）·“诸葛亮南征，平定益州郡、永昌等地叛乱并重置地方统治关系”的前期动员与资源准备}账本记录蜀汉与南中地方集团在建兴三年（225年）发生““诸葛亮南征，平定益州郡、永昌等地叛乱并重置地方统治关系”的前期动员与资源准备”。条目把背景概括为当时的权力、资源与信息约束推动相关过程展开，并指出这一过程为“诸葛亮南征，平定益州郡、永昌等地叛乱并重置地方统治关系”提供可观察的条件。这个节点进入区域社会时，要经由文书、道路、军队、市场、寺院和家庭等不同中介；因此，政治行动的宣布、地方接收和日常生活的改变不能被视作同一时刻完成。

本条依据《三国志·诸葛亮传》；《三国志·马谡传》裴松之注引《襄阳记》；《华阳国志·南中志》，证据提示为“核心事项已有既有账本来源；本分解节点的参与者、执行范围和时间先后仍须逐案核验”。这些材料可支持年序、官方决定、战事或特定地点的线索，却难以直接统计所有人的财富、迁徙、信仰和动机。更稳妥的读法是区分诏令与实施、军报与人口、行政称谓与自我认同，并让地方材料的缺环限制解释的范围。

由此可见，该事件不是孤立的年表点，而是资源、信息与权力关系重新配置的一环。不同地区的官员、社群和家庭可能以不同方式承担征发、保护、迁移或宗教活动；现存来源只能照亮其中一部分，不能替沉默者制造统一的经验。

\eventanchor{TK-LEDGER-009-2}{建兴三年（225年）·“诸葛亮南征，平定益州郡、永昌等地叛乱并重置地方统治关系”的命令、联盟或地方响应}账本记录蜀汉与南中地方集团在建兴三年（225年）发生““诸葛亮南征，平定益州郡、永昌等地叛乱并重置地方统治关系”的命令、联盟或地方响应”。条目把背景概括为当时的权力、资源与信息约束推动相关过程展开，并指出这一过程为“诸葛亮南征，平定益州郡、永昌等地叛乱并重置地方统治关系”提供可观察的条件。这个节点进入区域社会时，要经由文书、道路、军队、市场、寺院和家庭等不同中介；因此，政治行动的宣布、地方接收和日常生活的改变不能被视作同一时刻完成。

本条依据《三国志·诸葛亮传》；《三国志·马谡传》裴松之注引《襄阳记》；《华阳国志·南中志》，证据提示为“核心事项已有既有账本来源；本分解节点的参与者、执行范围和时间先后仍须逐案核验”。这些材料可支持年序、官方决定、战事或特定地点的线索，却难以直接统计所有人的财富、迁徙、信仰和动机。更稳妥的读法是区分诏令与实施、军报与人口、行政称谓与自我认同，并让地方材料的缺环限制解释的范围。

由此可见，该事件不是孤立的年表点，而是资源、信息与权力关系重新配置的一环。不同地区的官员、社群和家庭可能以不同方式承担征发、保护、迁移或宗教活动；现存来源只能照亮其中一部分，不能替沉默者制造统一的经验。

\eventanchor{TK-LEDGER-009-3}{建兴三年（225年）·诸葛亮南征，平定益州郡、永昌等地叛乱并重置地方统治关系}账本记录蜀汉与南中地方集团在建兴三年（225年）发生“诸葛亮南征，平定益州郡、永昌等地叛乱并重置地方统治关系”。条目把背景概括为“诸葛亮南征，平定益州郡、永昌等地叛乱并重置地方统治关系”发生的政治与区域条件共同作用，并指出该节点改变后续资源、联盟或制度处置的条件。这个节点进入区域社会时，要经由文书、道路、军队、市场、寺院和家庭等不同中介；因此，政治行动的宣布、地方接收和日常生活的改变不能被视作同一时刻完成。

本条依据《三国志·诸葛亮传》；《三国志·马谡传》裴松之注引《襄阳记》；《华阳国志·南中志》，证据提示为“核心事项已有既有账本来源；本分解节点的参与者、执行范围和时间先后仍须逐案核验”。这些材料可支持年序、官方决定、战事或特定地点的线索，却难以直接统计所有人的财富、迁徙、信仰和动机。更稳妥的读法是区分诏令与实施、军报与人口、行政称谓与自我认同，并让地方材料的缺环限制解释的范围。

由此可见，该事件不是孤立的年表点，而是资源、信息与权力关系重新配置的一环。不同地区的官员、社群和家庭可能以不同方式承担征发、保护、迁移或宗教活动；现存来源只能照亮其中一部分，不能替沉默者制造统一的经验。

\eventanchor{TK-LEDGER-009-4}{建兴三年（225年）·“诸葛亮南征，平定益州郡、永昌等地叛乱并重置地方统治关系”的执行中的空间与人员处置}账本记录蜀汉与南中地方集团在建兴三年（225年）发生““诸葛亮南征，平定益州郡、永昌等地叛乱并重置地方统治关系”的执行中的空间与人员处置”。条目把背景概括为“诸葛亮南征，平定益州郡、永昌等地叛乱并重置地方统治关系”发生的政治与区域条件共同作用，并指出该节点改变后续资源、联盟或制度处置的条件。这个节点进入区域社会时，要经由文书、道路、军队、市场、寺院和家庭等不同中介；因此，政治行动的宣布、地方接收和日常生活的改变不能被视作同一时刻完成。

本条依据《三国志·诸葛亮传》；《三国志·马谡传》裴松之注引《襄阳记》；《华阳国志·南中志》，证据提示为“核心事项已有既有账本来源；本分解节点的参与者、执行范围和时间先后仍须逐案核验”。这些材料可支持年序、官方决定、战事或特定地点的线索，却难以直接统计所有人的财富、迁徙、信仰和动机。更稳妥的读法是区分诏令与实施、军报与人口、行政称谓与自我认同，并让地方材料的缺环限制解释的范围。

由此可见，该事件不是孤立的年表点，而是资源、信息与权力关系重新配置的一环。不同地区的官员、社群和家庭可能以不同方式承担征发、保护、迁移或宗教活动；现存来源只能照亮其中一部分，不能替沉默者制造统一的经验。

\eventanchor{TK-LEDGER-009-5}{建兴三年（225年）·“诸葛亮南征，平定益州郡、永昌等地叛乱并重置地方统治关系”的后续制度或军政调整}账本记录蜀汉与南中地方集团在建兴三年（225年）发生““诸葛亮南征，平定益州郡、永昌等地叛乱并重置地方统治关系”的后续制度或军政调整”。条目把背景概括为核心节点发生后仍须处理其遗留问题，并指出后续影响须在不同地区和材料类型中分别核验。这个节点进入区域社会时，要经由文书、道路、军队、市场、寺院和家庭等不同中介；因此，政治行动的宣布、地方接收和日常生活的改变不能被视作同一时刻完成。

本条依据《三国志·诸葛亮传》；《三国志·马谡传》裴松之注引《襄阳记》；《华阳国志·南中志》，证据提示为“核心事项已有既有账本来源；本分解节点的参与者、执行范围和时间先后仍须逐案核验”。这些材料可支持年序、官方决定、战事或特定地点的线索，却难以直接统计所有人的财富、迁徙、信仰和动机。更稳妥的读法是区分诏令与实施、军报与人口、行政称谓与自我认同，并让地方材料的缺环限制解释的范围。

由此可见，该事件不是孤立的年表点，而是资源、信息与权力关系重新配置的一环。不同地区的官员、社群和家庭可能以不同方式承担征发、保护、迁移或宗教活动；现存来源只能照亮其中一部分，不能替沉默者制造统一的经验。

\eventanchor{TK-LEDGER-009-6}{建兴三年（225年）·“诸葛亮南征，平定益州郡、永昌等地叛乱并重置地方统治关系”的异源材料与影响范围的复核}账本记录蜀汉与南中地方集团在建兴三年（225年）发生““诸葛亮南征，平定益州郡、永昌等地叛乱并重置地方统治关系”的异源材料与影响范围的复核”。条目把背景概括为核心节点发生后仍须处理其遗留问题，并指出后续影响须在不同地区和材料类型中分别核验。这个节点进入区域社会时，要经由文书、道路、军队、市场、寺院和家庭等不同中介；因此，政治行动的宣布、地方接收和日常生活的改变不能被视作同一时刻完成。

本条依据《三国志·诸葛亮传》；《三国志·马谡传》裴松之注引《襄阳记》；《华阳国志·南中志》，证据提示为“核心事项已有既有账本来源；本分解节点的参与者、执行范围和时间先后仍须逐案核验”。这些材料可支持年序、官方决定、战事或特定地点的线索，却难以直接统计所有人的财富、迁徙、信仰和动机。更稳妥的读法是区分诏令与实施、军报与人口、行政称谓与自我认同，并让地方材料的缺环限制解释的范围。

由此可见，该事件不是孤立的年表点，而是资源、信息与权力关系重新配置的一环。不同地区的官员、社群和家庭可能以不同方式承担征发、保护、迁移或宗教活动；现存来源只能照亮其中一部分，不能替沉默者制造统一的经验。

\eventanchor{TK-LEDGER-010-1}{黄初七年五月（226年）·“曹丕去世，曹叡继位，是为魏明帝”的前期动员与资源准备}账本记录曹魏在黄初七年五月（226年）发生““曹丕去世，曹叡继位，是为魏明帝”的前期动员与资源准备”。条目把背景概括为当时的权力、资源与信息约束推动相关过程展开，并指出这一过程为“曹丕去世，曹叡继位，是为魏明帝”提供可观察的条件。这个节点进入区域社会时，要经由文书、道路、军队、市场、寺院和家庭等不同中介；因此，政治行动的宣布、地方接收和日常生活的改变不能被视作同一时刻完成。

本条依据《三国志·文帝纪》；《三国志·明帝纪》；《资治通鉴》魏纪，证据提示为“核心事项已有既有账本来源；本分解节点的参与者、执行范围和时间先后仍须逐案核验”。这些材料可支持年序、官方决定、战事或特定地点的线索，却难以直接统计所有人的财富、迁徙、信仰和动机。更稳妥的读法是区分诏令与实施、军报与人口、行政称谓与自我认同，并让地方材料的缺环限制解释的范围。

由此可见，该事件不是孤立的年表点，而是资源、信息与权力关系重新配置的一环。不同地区的官员、社群和家庭可能以不同方式承担征发、保护、迁移或宗教活动；现存来源只能照亮其中一部分，不能替沉默者制造统一的经验。

\eventanchor{TK-LEDGER-010-2}{黄初七年五月（226年）·“曹丕去世，曹叡继位，是为魏明帝”的命令、联盟或地方响应}账本记录曹魏在黄初七年五月（226年）发生““曹丕去世，曹叡继位，是为魏明帝”的命令、联盟或地方响应”。条目把背景概括为当时的权力、资源与信息约束推动相关过程展开，并指出这一过程为“曹丕去世，曹叡继位，是为魏明帝”提供可观察的条件。这个节点进入区域社会时，要经由文书、道路、军队、市场、寺院和家庭等不同中介；因此，政治行动的宣布、地方接收和日常生活的改变不能被视作同一时刻完成。

本条依据《三国志·文帝纪》；《三国志·明帝纪》；《资治通鉴》魏纪，证据提示为“核心事项已有既有账本来源；本分解节点的参与者、执行范围和时间先后仍须逐案核验”。这些材料可支持年序、官方决定、战事或特定地点的线索，却难以直接统计所有人的财富、迁徙、信仰和动机。更稳妥的读法是区分诏令与实施、军报与人口、行政称谓与自我认同，并让地方材料的缺环限制解释的范围。

由此可见，该事件不是孤立的年表点，而是资源、信息与权力关系重新配置的一环。不同地区的官员、社群和家庭可能以不同方式承担征发、保护、迁移或宗教活动；现存来源只能照亮其中一部分，不能替沉默者制造统一的经验。

\eventanchor{TK-LEDGER-010-3}{黄初七年五月（226年）·曹丕去世，曹叡继位，是为魏明帝}账本记录曹魏在黄初七年五月（226年）发生“曹丕去世，曹叡继位，是为魏明帝”。条目把背景概括为“曹丕去世，曹叡继位，是为魏明帝”发生的政治与区域条件共同作用，并指出该节点改变后续资源、联盟或制度处置的条件。这个节点进入区域社会时，要经由文书、道路、军队、市场、寺院和家庭等不同中介；因此，政治行动的宣布、地方接收和日常生活的改变不能被视作同一时刻完成。

本条依据《三国志·文帝纪》；《三国志·明帝纪》；《资治通鉴》魏纪，证据提示为“核心事项已有既有账本来源；本分解节点的参与者、执行范围和时间先后仍须逐案核验”。这些材料可支持年序、官方决定、战事或特定地点的线索，却难以直接统计所有人的财富、迁徙、信仰和动机。更稳妥的读法是区分诏令与实施、军报与人口、行政称谓与自我认同，并让地方材料的缺环限制解释的范围。

由此可见，该事件不是孤立的年表点，而是资源、信息与权力关系重新配置的一环。不同地区的官员、社群和家庭可能以不同方式承担征发、保护、迁移或宗教活动；现存来源只能照亮其中一部分，不能替沉默者制造统一的经验。

\eventanchor{TK-LEDGER-010-4}{黄初七年五月（226年）·“曹丕去世，曹叡继位，是为魏明帝”的执行中的空间与人员处置}账本记录曹魏在黄初七年五月（226年）发生““曹丕去世，曹叡继位，是为魏明帝”的执行中的空间与人员处置”。条目把背景概括为“曹丕去世，曹叡继位，是为魏明帝”发生的政治与区域条件共同作用，并指出该节点改变后续资源、联盟或制度处置的条件。这个节点进入区域社会时，要经由文书、道路、军队、市场、寺院和家庭等不同中介；因此，政治行动的宣布、地方接收和日常生活的改变不能被视作同一时刻完成。

本条依据《三国志·文帝纪》；《三国志·明帝纪》；《资治通鉴》魏纪，证据提示为“核心事项已有既有账本来源；本分解节点的参与者、执行范围和时间先后仍须逐案核验”。这些材料可支持年序、官方决定、战事或特定地点的线索，却难以直接统计所有人的财富、迁徙、信仰和动机。更稳妥的读法是区分诏令与实施、军报与人口、行政称谓与自我认同，并让地方材料的缺环限制解释的范围。

由此可见，该事件不是孤立的年表点，而是资源、信息与权力关系重新配置的一环。不同地区的官员、社群和家庭可能以不同方式承担征发、保护、迁移或宗教活动；现存来源只能照亮其中一部分，不能替沉默者制造统一的经验。
